\documentclass[11pt]{article}
\usepackage[letterpaper,margin=0.82in]{geometry}
\usepackage[T1]{fontenc}
\usepackage{lmodern,microtype,amsmath,amssymb,amsthm,mathtools}
\usepackage{booktabs,longtable,array,enumitem,xcolor,fancyhdr}
\usepackage{hyperref}
\hypersetup{colorlinks=true,linkcolor=blue!45!black,urlcolor=blue!55!black,citecolor=blue!45!black,
 pdftitle={A fixed-base modular Diophantine equation for ZFC consistency},
 pdfauthor={Research prepared with AI assistance for Vladimir Reshetnikov}}
\setlength{\headheight}{14pt}
\pagestyle{fancy}\fancyhf{}
\fancyhead[L]{\small Diophantine consistency: fixed-base and modular certificates}
\fancyfoot[C]{\thepage}
\setlength{\parindent}{0pt}\setlength{\parskip}{5pt plus 1pt}
\setcounter{tocdepth}{1}
\setlist{nosep,leftmargin=*}
\emergencystretch=3em
\newtheorem{theorem}{Theorem}[section]
\newtheorem{lemma}[theorem]{Lemma}
\newtheorem{proposition}[theorem]{Proposition}
\newtheorem{corollary}[theorem]{Corollary}
\theoremstyle{remark}\newtheorem{remark}[theorem]{Remark}
\newcommand{\Z}{\mathbb Z}\newcommand{\N}{\mathbb N}
\newcommand{\Con}{\operatorname{Con}}
\newcommand{\lcm}{\operatorname{lcm}}
\newcommand{\G}{\mathcal G}\newcommand{\LL}{\mathcal L}
\newcommand{\abs}[1]{\lvert #1\rvert}
% BEGIN report/metrics.tex
\newcommand{\MainVars}{230}
\newcommand{\MainRows}{176}
\newcommand{\MainTerms}{4,482}
\newcommand{\MainHeight}{7,968}
\newcommand{\MainLastRow}{175}
\newcommand{\MainLastVar}{229}
\newcommand{\CompactVars}{229}
\newcommand{\CompactRows}{175}
\newcommand{\CompactTerms}{12,006}
\newcommand{\CompactLastRow}{174}
\newcommand{\CompactLastVar}{228}
% END report/metrics.tex
\begin{document}
\begin{titlepage}
\vspace*{0.8in}
{\color{blue!45!black}\large RESEARCH CONTRIBUTION FOR PROVEIT}\par
\vspace{0.22in}
{\huge\bfseries A fixed-base, modular Diophantine equation for ZFC consistency\par}
\vspace{0.25in}
{\Large Reusing a binomial modulus as the proof-label base\par}
\vspace{0.35in}
{\large Prepared with AI assistance for Vladimir Reshetnikov\par}
{\large 25 October 2026\par}
\vspace{0.35in}
\begin{center}
\fbox{\begin{minipage}{0.86\textwidth}
\textbf{Main delivered equation.} An explicitly specified polynomial in
\textbf{\MainVars\ unrestricted integer variables}, of exact degree four,
defined by \MainRows\ quadratic residuals. Its collected expansion has
\MainTerms\ nonzero monomials and coefficient height \MainHeight.
\[
 P(\boldsymbol z)=\sum_{i=0}^{\MainLastRow}R_i(\boldsymbol z)^2=0.
\]
The ordinary proof, composed with the retained numerical-calculus
calibration, establishes
\[
 \mathrm{ZFC}\vdash
 \bigl(\neg\exists\boldsymbol z\in\Z^{\MainVars}\ P(\boldsymbol z)=0\bigr)
 \ \longleftrightarrow\ \Con(\mathrm{ZFC}).
\]
Every residual is printed in Appendix~\ref{app:residuals} and supplied
as exact machine-readable data. No unspecified coefficient or variable
exponent remains in the delivered polynomial.
\end{minipage}}
\end{center}
\vfill
\textbf{Evidence boundary.} This is an ordinary mathematical proof with
reproducible finite algebraic checks. It is not a new end-to-end Lean/Rocq
formalization, a proof of consistency, or a claim of global minimality.
\end{titlepage}

\begin{abstract}
We simplify the direct proof-certificate route to a Diophantine equation
for syntactic consistency of ZFC. The central change is to give one
binomial coefficient two jobs. With a proof horizon $t$, set
$R=16^t$ and $B=\binom{R+t}{t}$. The number $B$ supplies both the
prime-power selectors for proof rows and the evaluation point in the
label weights $1+jB$. An elementary bound for $\lcm(1,\ldots,t)$
proves the required row coverage, while $B>(t+1)^t$ proves exact
label-multiplicity recovery. No certificate of factorial authenticity,
and no separate variable-base growth power, is needed.

We then compose this arithmetic front with the predecessor's modular
numerator method in the current signed Pell architecture. The new
compatibility argument allows all nine CRT representatives to remain
signed, recovers the necessary positive quantities from existing power
requests, and moves eleven numerator outputs out of the magnitude frame.
A new cofinal two-square completion of the Pell center restores the
common modular-root congruence. Only six exponential requests remain exact; successor
denominators use multiplication. The resulting finite quadratic system
is contracted with explicit integral restoration maps and converted to
the literal quartic displayed in this report. The contribution is a
structural improvement within the direct, small-coefficient family;
few-variable constructions with different coefficient and degree costs
remain distinct alternatives.
\end{abstract}
\tableofcontents
\newpage

\section{The exact result and its scope}

We use $\N=\{0,1,2,\ldots\}$ and quantify every coordinate of the
principal equation over all of $\Z$. The file
\nolinkurl{equations/balanced/equation.json} declares the ordered variables
and every quadratic residual. A sparse term $[c,[i_1,\ldots,i_d]]$ denotes
$c z_{i_1}\cdots z_{i_d}$; an empty index list denotes the constant $c$.
The finite sum of residual squares is the definition of $P$. The separate
\nolinkurl{expanded_polynomial.json} gives every collected coefficient.

\begin{theorem}[Consistency equivalence]\label{thm:main}
For the supplied fixed polynomial $P\in\Z[z_0,\ldots,z_{\MainLastVar}]$,
ZFC proves that an integer zero exists if and only if a finite syntactic
ZFC refutation exists. Consequently ZFC proves
\[
 \neg\exists\boldsymbol z\in\Z^{\MainVars}\ P(\boldsymbol z)=0
 \quad\longleftrightarrow\quad\Con(\mathrm{ZFC}).
\]
The exact degree of $P$ is four.
\end{theorem}

The companion in \nolinkurl{equations/compact/equation.json} uses
\CompactVars\ integer variables and \CompactRows\ quadratic residuals.
It has \CompactTerms\ collected monomials and the same coefficient
height \MainHeight. It is obtained by one additional integral affine
elimination, so it has an integer zero exactly when $P$ does.
Appendix~\ref{app:compact} prints its full residual list as well.
The principal version is preferred when the shorter expansion matters;
the companion is preferred when one fewer coordinate matters more.

The proof does not assume either consistency or inconsistency. In
particular, proving the equivalence does not prove the absence of zeros.
The syntactic consistency predicate is the one calibrated by the
repository's fixed numerical calculus $C_{222}$, with truth label $0$
and falsity label $1$. The retained logical interfaces are
\nolinkurl{KEndpoint.c222_iff_zf_refutation} and
\nolinkurl{ConstructibleChoice.zf_refutation_iff_zfc_refutation}.
Their source and provenance accompany the package; their complete
transitive proof-assistant development was not rebuilt here.

The new arithmetic arguments below are finite number theory and
primitive recursive coding arguments formalizable in ZFC. Their
composition with these inherited calibrations proves
Theorem~\ref{thm:main}. A fresh machine-checked object-language ZFC
derivation is a separate, unfinished task. Exact coefficient checks
verify the finite artifacts and their algebraic transformations, not
the unbounded semantic theorem by themselves.

\subsection{What ``simpler'' means in this report}

There is no single total ordering of variable count, degree, coefficient
height, monomial count, arithmetic circuit cost, and proof length. The
repository already contains equations with far fewer unknowns, including
quartics with large fixed coefficient payloads. This report optimizes the
direct proof-calculus family, which retains modest literal coefficients
and avoids both a $K215$ simulation and an instantiation of Jones's
universal equation.

The comparison snapshot is the immutable ProveIt commit
\begin{center}\small
\texttt{b7e35224c04fe1031c0a53f80a395e9b32205022}.
\end{center}
We inspected its active research documents, delivered packages through
147, and all eight incoming archives present at that commit. The
top-level comparison tables lagged behind the newest incoming material;
the comparison below therefore uses the actual incoming endpoints.

% BEGIN report/comparison.tex
\begin{center}\small
\begin{tabular}{@{}lrrrrr@{}}
\toprule
Profile & Variables & Rows & Degree & Monomials & Height \\
\midrule
Matched parent & 252 & 182 & 4 & 5,340 & 128,446 \\
Factorial-free, sparse & 242 & 178 & 4 & 5,018 & 7,968 \\
Factorial-free, compact & 241 & 177 & 4 & 12,774 & 7,968 \\
Shared error, minimum & 240 & 177 & 4 & 21,482 & 6,144 \\
Shared error, lower height & 241 & 178 & 4 & 9,544 & 4,512 \\
Normalized, balanced & 244 & 176 & 4 & 5,212 & 3,968 \\
\midrule
\textbf{New principal} & \textbf{230} & \textbf{176} & \textbf{4} & \textbf{4,482} & \textbf{7,968} \\
\textbf{New compact} & \textbf{229} & \textbf{175} & \textbf{4} & \textbf{12,006} & \textbf{7,968} \\
\bottomrule
\end{tabular}
\end{center}
Every row in this table quantifies over unrestricted integers.
Height means the largest absolute coefficient after collecting the expanded quartic.
The source paths and exact metric files for the comparison rows are in \nolinkurl{provenance/comparison.json}.
% END report/comparison.tex

The modular numerator device itself is inherited: the retained
\nolinkurl{INHERITED_MODULAR_FACTORIAL_PROOF.md} already uses six exact
powers, eleven modular numerators, and a common Pell-root congruence.
What is changed here is its compatibility with the fixed-base/binomial
weight front and the later signed, pooled certificate architecture.
The dual use of $B$, the new sign bootstrap, the two-square cofinal-center
theorem, and their complete composition were not found in the inspected
snapshot. No literature-wide priority claim is made.

\section{The arithmetic design before polynomial compilation}

The following notation explains certified register values. Powers and
binomial coefficients in this section are not extra operations in the
final quartic. They are replaced by the polynomial certificates discussed
below and printed literally in the appendix.

Choose a cofinal proof horizon $t=16+h^2$. The essential front is
\begin{align}
 R&=16^t,& B&=2b+1,& B&=\binom{R+t}{t},\label{eq:front1}\\
 y&=RB(2v+1)^2,&
 H&=2+2(R+t)+C^2+\sum_{j=1}^{8}A_j^2,\label{eq:front2}\\
 N&=B^x,& q&=(BH)^x=6z+3.&\label{eq:front3}
\end{align}
Here $C,A_1,\ldots,A_8$ are the actual CRT representatives. Some raw
stored coordinates equal these values minus one; every literal source
edit preserves that distinction. The nine representatives are unrestricted
integers. Their small residues recovered later are nonnegative.

The two beta endpoints and the row-index congruence are
\begin{align}
 x&=B+1+(y+1)u_0,\label{eq:initial}\\
 x&=1+[1+(t+1)y]u_t,\label{eq:final}\\
 B h_B&=C+R+1.\label{eq:index}
\end{align}
Every Horn weight is now $1+jB$. The former independent growth base
$F=[2(t+1)]^t$ and the radius $F^2$ are removed from the semantics.
The same $B$ is used throughout; changing only some occurrences would
not give the claimed construction.

The six exact requests are
\begin{equation}\label{eq:six}
 R=16^t,\quad N=B^x,\quad q=(BH)^x,\quad
 D_t=q^t,\quad D_y=q^{y+1},\quad M=q^x.
\end{equation}
Set $D_{t+1}=qD_t$ by multiplication. The eleven binomial numerator
indices are $R+t$, $C$, the eight representatives $A_j$, and the
derived label $J(\boldsymbol A)$. Their relevant lower indices are
$t,t+1,y+1,x$. The exact-binomial digit for $B$ uses $t$;
the index range uses $t$; four syntax ranges and the derived-label
range use $t+1$; two accumulator ranges use $y+1$; two quotient
ranges use $x$.

\section{A fixed-base radius and a binomial with two roles}

\subsection{The radius supplies every row selector}

Let $\Lambda_t=\lcm(1,\ldots,t)$, with $t\ge1$.

\begin{lemma}\label{lem:lcm}
For every $n\ge1$, $\Lambda_n\le4^{n-1}$. In particular,
$16^t\ge\Lambda_t^2$.
\end{lemma}
\begin{proof}
Let $m=\lceil n/2\rceil$. The quotient $\Lambda_n/\Lambda_m$
divides $\binom{n}{\lfloor n/2\rfloor}$. To see this, its valuation
at any prime is at most one: there is at most one new prime power
between $m+1$ and $n$. If such a prime power occurs, its contribution
to the factorial valuation of the central binomial coefficient is one,
and all other contributions are nonnegative.

For $n=2r$, induction and the binomial row sum give
\[
 \Lambda_{2r}\le\Lambda_r\binom{2r}{r}
 \le4^{r-1}2^{2r}=4^{2r-1}.
\]
For $n=2r+1$, the two equal central coefficients together are at most
the row sum, so $\binom{2r+1}{r}\le2^{2r}$. Hence
\[
 \Lambda_{2r+1}\le\Lambda_{r+1}\binom{2r+1}{r}
 \le4^r4^r=4^{2r}.
\]
The initial case $n=1$ is immediate. The stated radius consequence
is weaker than the resulting strict bound and follows at once.
\end{proof}

The next selector theorem is inherited from the incoming factorial-free
binomial CRT report, and is restated with its proof because it is the
key interface used by the new fixed radius. The classical lcm/binomial
divisibility used in its proof is also related to Farhi's identity
\cite{farhi}.

\begin{lemma}[Coverage and compatibility of selectors]\label{lem:selectors}
Let $t\ge1$, $R\ge\Lambda_t^2$, and $B=\binom{R+t}{t}$.
For every $i\in\{1,\ldots,t\}$ there is a full prime power
$p^e\Vert B$ such that $p^e\mid R+i$ and $p^e>t$.
Conversely, every full prime power $p^e\Vert B$ divides at least one
of $R+1,\ldots,R+t$.
\end{lemma}
\begin{proof}
For the first assertion, $R+i>\Lambda_t^2$, so some prime $p$ has
$v=v_p(R+i)>2a$, where $a=\lfloor\log_p t\rfloor$. Since $p^v>t$,
the valuation of $R+j$ for $j\ne i$ is exactly $v_p(j-i)$. Therefore
\[
 e=v_p(B)=v-\delta,\qquad
 \delta=v_p\!\left(t\binom{t-1}{i-1}\right).
\]
For each $p^r$, the corresponding factorial-valuation summand for
$\delta$ is zero or one, and vanishes when $p^r>t$.
Thus $0\le\delta\le a$, so $a<e\le v$ as required.

For the converse, put $v=\max_i v_p(R+i)$. The valuation of $B$ is
the sum over $r\ge1$ of
\[
 \left\lfloor\frac{R+t}{p^r}\right\rfloor
 -\left\lfloor\frac R{p^r}\right\rfloor
 -\left\lfloor\frac t{p^r}\right\rfloor.
\]
Each term is zero or one and all terms above $r=v$ vanish. Hence
$e\le v$, and a factor attaining the maximum is compatible.
\end{proof}

The two directions have different uses. Coverage proves soundness for
each row. Compatibility supplies assignments at every prime of the
global modulus in the CRT completeness construction, including primes
not chosen as a particular row's selector.

\subsection{The same binomial separates label multiplicities}

\begin{lemma}\label{lem:binomial-growth}
For $t\ge1$ and $R=16^t$,
\[
 B=\binom{R+t}{t}>(t+1)^t.
\]
Moreover $B$ is odd.
\end{lemma}
\begin{proof}
The inequality $16^t>t(t+1)$ follows by induction: its left side
grows by a factor of 16 and the right side by at most 3. Every factor
$(R+i)/i$, $1\le i\le t$, is greater than $t+1$, proving the first
claim by the product formula for $B$.

For oddness, $v_2(R)=4t>\lfloor\log_2t\rfloor$. Consequently
$v_2(R+i)=v_2(i)$ for every $1\le i\le t$. The valuations in the
numerator product and in $t!$ cancel exactly.
\end{proof}

Here is the large-evaluation argument inherited from the factorial-free
predecessor, now applied to $B$. Suppose the decoded $t$ proof rows
have conclusion labels $c_i$ and premise labels $p_i,q_i$ in $[0,t]$,
and the accumulator updates imply
\begin{equation}\label{eq:balance}
 \prod_{i<t}(1+c_iB)
 =(1+B)\prod_{i<t}(1+p_iB)(1+q_iB).
\end{equation}
Discard factors with label zero. If the two sides have different numbers
of remaining factors, the smaller factor count is at most $t$, while
$B>(t+1)^t$ makes a product with one more factor strictly larger than
any product with the smaller count. Thus the counts agree and are at
most $t$.

Replace $B$ by an indeterminate $T$. Each resulting product has
nonnegative coefficients and coefficient sum at most $(t+1)^t<B$.
If their difference were nonzero, let $a_rT^r$ be its lowest nonzero
term. The equality at $T=B$ would imply $B\mid a_r$, although
$0<\abs{a_r}<B$. Thus the products agree in $\Z[T]$. Unique
factorization, or their distinct rational roots $-1/j$, gives the
exact multiplicity balance for each positive label.

To finish the flow argument, let $D$ be the set of labels derivable by
the valid Horn rules. It contains truth. A rule whose conclusion is
outside $D$ has at least one premise outside $D$. Summing conclusion
counts minus premise counts over the complement of $D$ therefore gives
a nonpositive number. If falsity were outside $D$, the exact balance
\eqref{eq:balance} would give one instead. Hence falsity is derivable.
No acyclicity assumption about the chosen multiset of rules is needed.

\section{Pell values, the free center, and modular numerators}

For an integer $A$, define the Pell polynomials by
\[
 X_0=1,\quad Y_0=0,\qquad
 X_{k+1}=AX_k+(A^2-1)Y_k,\quad Y_{k+1}=X_k+AY_k,
\]
and extend them to negative indices by $X_{-k}=X_k$ and $Y_{-k}=-Y_k$.
The retained square frame, semantic anchor, and pooled error certificate
recover exactly these pairs at every requested signed index. The
complete fourteen-pair orientation pool is retained.

For a genuine exact request with base $n$, the certified congruence is
\[
 X_k(A)+(n-A)Y_k(A)\equiv m
 \pmod{2An-n^2-1}.
\]
The magnitude frame bounds every exact base, exponent, and output atom
in absolute value by $W$. Its large-center bound makes the modulus
larger than the possible error, giving $m=n^k$ for $\abs n\ge2$ and
forcing $k\ge0$. If $k<0$, the same recurrence would force the
impossible integer equality $m n^{-k}=1$. These signed recovery
arguments are inherited from the retained source proofs.

\subsection{Why the center must be freed}

The current signed frame restricts its abscissa to a particular Pell
sequence. This alone does not justify choosing it in an arbitrary
congruence class. Retain the lower frame with abscissa $L$ and use
\begin{equation}\label{eq:free-center}
                         A=L^2+u^2+v^2.
\end{equation}
This gives $A\ge L^2$ at every integer solution. Completeness requires
a special argument: an arbitrary prescribed nonnegative difference is
not a sum of two squares. The frame itself is existentially chosen,
and its finite residue behavior provides the missing freedom. The
cofinal completion lemma below proves exactly what is needed. It saves
one coordinate relative to the initially verified constructor
$A=L^2+u^2+v^2+w(w+1)$.

Only the lower frame uses $L$; the principal Pell pairs and the
semantic anchor use $A$. The magnitude frame excludes the eleven
modular numerator outputs, which now depend on $A$. It also removes
redundant atoms with separately proved bounds. In particular, keeping
$t$, $R+t$, and $q$ controls the fixed base 16, $R$, $t+1$, and $q+1$:
$W\ge t^2+2\ge258$ and
\[
 R^2\le2\bigl((R+t)^2+t^2\bigr)\le2(W-2)<W^2.
\]
The successor output $D_{t+1}$ no longer needs an exact-power magnitude
bound. Every retained exact request and modular index still has its
required strict bound. These changes are part of the literal frame
certificate, rather than a numerical estimate on sample inputs.

\subsection{A single modular-root equation}

Set $\Omega=NM$ and impose
\begin{equation}\label{eq:common-root}
 2A(q+1)-(q+1)^2-1=\Omega w.
\end{equation}
For every signed integer index $n$, retain the polynomial numerator
\begin{equation}\label{eq:mu}
 \mu_n=X_n(A)+(q+1-A)Y_n(A).
\end{equation}
It need not equal the enormous integer $(q+1)^n$.

\begin{lemma}[Signed modular powers]\label{lem:modular}
Under \eqref{eq:common-root}, with $\Omega\ne0$,
\[
 \mu_n\equiv(q+1)^n\pmod\Omega\qquad(n\in\Z),
\]
where a negative exponent means the inverse power in the residue ring.
\end{lemma}
\begin{proof}
Write $s=q+1$. Equation~\eqref{eq:common-root} says
$s(2A-s)\equiv1\pmod\Omega$, so $s$ is a unit. The Pell
recurrence and $(s-A)^2\equiv A^2-1$ show that the evaluated Pell
unit at index $n$ is $s^n$. This proves positive indices by induction,
and inverse powers give negative indices.
\end{proof}

In completeness, $B$ and $q$ are odd, and $N\mid q$; hence
$\Omega$ is odd and every prime dividing $\Omega$ divides $q$.
The root congruence specifies
\[
 A_0=\frac{(q+1)^2+1}{2(q+1)}
    =1+\frac{q^2}{2(q+1)}\pmod\Omega.
\]
Its denominator is invertible. The following lemma constructs this
residue class using precisely the two-square center.

\begin{theorem}[Cofinal two-square center]\label{thm:two-square}
Let $W\ge2$ and let $q,\Omega$ be positive odd integers such that
every prime of $\Omega$ divides $q$. There exist a positive odd $k$
and integers $u,v$ such that, with $L=X_{Wk}(W+1)$,
\[
 A=L^2+u^2+v^2\equiv A_0\pmod\Omega.
\]
The value $A$ can be made arbitrarily large, and $Y_{Wk}(W+1)/W$
is an odd integer. Thus the old odd-quotient frame and the new common
root remain simultaneously complete.
\end{theorem}
\begin{proof}
If $\Omega=1$, choose $k=1$ and arbitrary sufficiently large $u$;
the frame-parity argument below still applies. Assume $\Omega>1$.
Put $D=W(W+2)$ and work in the finite quadratic ring
\[
 S_\Omega=(\Z/\Omega\Z)[s]/(s^2-D).
\]
The element $e=(W+1)+s$ is a norm-one unit. Write the finite order of
$e^W$ as $2^r k$ with $k$ odd. Then $g=e^{Wk}$ has power-of-two
order. Its coefficients are $L=X_{Wk}(W+1)$ and $Y=Y_{Wk}(W+1)$.

Fix $p^a\Vert\Omega$. If $L\equiv\varepsilon\pmod p$ with
$\varepsilon\in\{1,-1\}$, then the image of $g$ is exactly
$\varepsilon$ in $S_p$. When $D\not\equiv0\pmod p$, this follows
from $L^2-DY^2=1$. When $D\equiv0\pmod p$, the nilpotent part of
$\varepsilon+Ys$ has $p$-power order, and is killed only trivially by
a power of two because $p$ is odd. This proves the assertion also
in the ramified case.

The kernel of the unit reduction map $S_{p^a}^{\times}\to S_p^{\times}$
is a finite $p$-group: its elements are $1+pz$. The element
$\varepsilon g$ belongs to this kernel and has power-of-two order,
so it is the identity. Consequently $L^2\equiv1\pmod{p^a}$.

Since $A_0\equiv1\pmod p$, the difference $A_0-L^2$ is a unit
unless $L\equiv\pm1\pmod p$. In the latter case it equals
$q^2/(2(q+1))$ modulo $p^a$, and is either zero or has the even
valuation $2v_p(q)$. Every residue of zero or even valuation modulo
an odd prime power is a sum of two squares. For a unit, the two
sets $\{a^2:a\in\mathbb F_p\}$ and
$\{c-b^2:b\in\mathbb F_p\}$ intersect for any nonzero residue $c$,
since each has $(p+1)/2$ elements. At least
one resulting coordinate is nonzero, so its invertible derivative
lifts the solution to every $p^a$. For valuation $2b$, multiply
both coordinates of a unit representation by $p^b$.

CRT, applied separately to the two square coordinates, therefore gives
$u^2+v^2\equiv A_0-L^2\pmod\Omega$. Replacing $u$ by
$u+j\Omega$ makes $A$ arbitrarily large.

It remains to check the literal frame parity. The recurrence modulo $W$
gives $W\mid Y_W(W+1)$. If $W$ is odd, the ordinate is odd.
If $W$ is even, its base is odd, every abscissa is odd, and
$Y_{2m}=2X_mY_m$ gives $v_2(Y_W)=v_2(W)$; odd-index ordinates are
odd. Hence $Y_W(W+1)/W$ is odd in both cases. Finally,
\[
 Y_{Wk}(W+1)=Y_W(W+1)Y_k(X_W(W+1)),
\]
whose second factor is odd for odd $k$. Thus the quotient remains
odd and supplies the old witness in
$L^2-W^3(W+2)(2z+1)^2=1$.
\end{proof}

\begin{remark}[Why the frame choice matters]
This theorem does not work by holding the least frame point fixed.
For example, $W=4$, $q=3$, $\Omega=9$ give the least frame abscissa
$L=4801$ and $A_0-L^2\equiv3\pmod9$, which is not a sum of two
squares modulo 9. Choosing the odd multiplier from the finite unit
order is essential. The proof does not add an order computation or
an extra quantified object to the polynomial; it chooses an existing
existential frame witness in the completeness direction.
\end{remark}

\section{Signed binomial digits from modular data}

For $n\in\Z$ and $k\in\N$, use the integer-valued generalized
binomial coefficient
\[
 \binom nk=\frac{n(n-1)\cdots(n-k+1)}{k!},\qquad \binom n0=1.
\]
For $n=-a<0$, this equals $(-1)^k\binom{a+k-1}{k}$.

\begin{lemma}[Uniform coefficient bound]\label{lem:signed-bound}
For every $n\in\Z$ and $k\ge0$,
\[
 \abs{\binom nk}\le\max(1,\abs n)^k.
\]
\end{lemma}
\begin{proof}
The ordinary nonnegative case is immediate. For $n=-a<0$,
$\binom{a+k-1}{k}$ counts multisets of size $k$ from $a$ colors.
Every multiset has a sorted word, and there are $a^k$ words. The
case $k=0$ is included.
\end{proof}

\begin{lemma}[The range implication for signed upper indices]\label{lem:signed-range}
If $p$ is prime, $e,k\ge1$, and $p^e\mid\binom nk$ for
$n\in\Z$, then $n\bmod p^e<k$.
\end{lemma}
\begin{proof}
If the falling factorial contains zero, then $0\le n<k$, and the
claim is immediate. Otherwise its valuation minus $v_p(k!)$ is a
sum of terms, each zero or one: a block of $k$ consecutive integers
contains either $\lfloor k/p^r\rfloor$ or one more multiples of
$p^r$. If the block contained no multiple of $p^e$, then $k<p^e$
and every term at level $r\ge e$ would vanish. The valuation would
be at most $e-1$, a contradiction. A multiple in the block gives
$n\equiv j\pmod{p^e}$ with $0\le j<k$.
\end{proof}

The implication is one-sided. Its converse is not used. Completeness
will use higher CRT precision to compensate for division by $k!$.

\subsection{Finite truncation, including negative indices}

The formal identity $(1+T)^n=\sum_{j\ge0}\binom njT^j$ holds
in $\Z[[T]]$ for all signed $n$, using the formal inverse for $n<0$.
It may equally be proved modulo $T^{k+1}$ by a finite recurrence,
so no analytic convergence claim is involved. Suppose $q>0$, $N\mid q$,
$1\le k\le x$, and $D=q^k$. Define
\[
 r_0=\sum_{j=0}^{k-1}\binom njq^j.
\]
The modulus $ND$ divides $\Omega=Nq^x$, and the omitted terms
of degree at least $k+1$ vanish modulo $ND$. Lemma~\ref{lem:modular}
therefore gives
\begin{equation}\label{eq:truncation}
 \mu_n\equiv r_0+D\binom nk\pmod{ND}.
\end{equation}
This remains true at the largest lower index $k=x$ because $N\mid q$.

\subsection{The common centered norm}

Write $\G(a,b,c)=a^2+b^2+c(c+1)$. This is nonnegative on
$\Z^3$, and represents every nonnegative integer. Indeed, Legendre's
three-square theorem represents $4g+1$ as three squares. Exactly one
square is odd modulo four; write the coordinates as $2a,2b,2c+1$.
Dividing the resulting equality by four gives
$g=a^2+b^2+c(c+1)$. This is the inherited ternary gap constructor;
it is distinct from the new two-square center theorem.

For each of the eleven digits, let $D_i$ be its already certified
denominator and put $D_i v_i=M$. Retain the normalized row
\begin{equation}\label{eq:normalized}
 v_i\mu_{n_i}=M Q_i+\rho_i.
\end{equation}
For the exact digit $Q_i=B+q h_i$; for a range digit $Q_i=N h_i$.
For $D_i=M$, $v_i=1$ is a constant rather than a new coordinate.
The single norm budget is
\begin{equation}\label{eq:norm}
 M^2=1+4\sum_{i=1}^{11}\rho_i^2+\G(\alpha,\beta,\gamma).
\end{equation}

Since $v_i\mid M$, equation~\eqref{eq:normalized} implies
$v_i\mid\rho_i$. Thus $r_i=\rho_i/v_i$ is an integer and
$\mu_{n_i}=D_iQ_i+r_i$. The norm gives $\abs{r_i}<D_i/2$.
If $\abs{r_0}<D_i/32$, then $D_i\mid r_i-r_0$ and
$\abs{r_i-r_0}<D_i$, so $r_i=r_0$. Applying
\eqref{eq:truncation} to a range digit yields
\[
                         N\mid\binom{n_i}{k_i}.
\]
For the exact digit, the stronger modulus $qD_t\mid\Omega$
yields $B\equiv\binom{R+t}{t}\pmod q$.

The argument does not require a signed representative's truncation to
be a least nonnegative remainder. Its small absolute value is sufficient.
Canonical completeness representatives are positive, but soundness
applies to every integer zero of the final equation.

\section{Soundness in a noncircular order}

\subsection{Recover positive arithmetic before reading digits}

The inherited Pell index and orientation proof applies first. Its
magnitude frame now contains the six exact requests and every modular
numerator index; it does not contain the modular outputs.
The free-center row gives $A>0$ and the required magnitude separation.

The request with base 16 gives $R=16^t>1$. The affine radix condition
$q=6z+3$ gives $\abs q\ge3$, and the explicitly odd $B$ is nonzero.
The exact request $D_y=q^{y+1}$ forces $y+1\ge0$. But
$y=RB(2v+1)^2$ is a nonzero multiple of $R$. A nonpositive $y$
would therefore be at most $-R<-1$. Thus $y>0$ and $B>0$.

The possible value $B=1$ needs separate treatment before binomial
authentication. For every signed $k$, the Pell evaluation at base 1
is congruent to 1 modulo $2(A-1)$: the Pell unit evaluates to the unit 1
in that residue ring. Hence the retained exact request forces
$N\equiv1\pmod{2(A-1)}$. Its frame bound $\abs N<W$ and the
large-center bound force $N=1$. Thus $N=B^x$ holds also in this
temporary base-one case. For $B\ge3$ it follows from the usual exact
power lemma once $x\ge0$ is recovered below.

The square-defined $H$ is positive for arbitrary signed representatives.
Hence the request $q=(BH)^x$ gives $x\ge0$ and $q>0$. The affine
radix excludes $x=0$. Since $y\ge RB>B$, a negative $u_0$ in
\eqref{eq:initial} would imply $x\le B-y<0$. Therefore
$u_0\ge0$ and $x\ge B+1>1$. Equation~\eqref{eq:final} now
forces $u_t\ge1$, giving
\begin{equation}\label{eq:large-x}
 x\ge2+(t+1)y,\qquad y\ge RB>B,\qquad x>2^{32}.
\end{equation}
This order uses neither the sign of a CRT representative nor the value
of a binomial digit. In particular, it avoids using factoriality or a
binomial interpretation to prove the signs needed for that interpretation.

All four denominator groups are now positive divisors of $M=q^x$,
and $N=B^x$ is positive and divides $q$. The local majorant and
derived-label budgets can be interpreted at this point.

\subsection{Bounds for every relevant coefficient}

Set $\ell=2^{22}$ and $E=2^{23}$. The retained local arithmetic
ledger bounds the derived-label polynomial in absolute value by
$H^\ell$ and bounds the decoded nonnegative local rule polynomial by
\[
 0\le\LL<(2+t+x+y+B)^E.
\]
The signed derived-label bound has a separate structural audit. For
$H\ge2$, assign exponent 1 to each bounded stored input. At an actual
sparse gate $g=\sum_m c_m\prod_{i\in m}v_i$, the auditor assigns
\[
 e(g)=\left\lceil\log_2\sum_m\abs{c_m}\right\rceil
       +\max_m\sum_{i\in m}e(v_i).
\]
Induction proves $\abs g\le H^{e(g)}$. The supplied ledger checks
91 ancestor gates and only four free leaves. It obtains exponent
$3{,}649{,}856$ for the stored wire $J-1$, and $3{,}649{,}857<\ell$
for $J$. It does not assume positivity of the inputs. The stored
inputs are actual representatives minus one; $H\ge2+a^2$ implies
$\abs{a-1}\le H$ for every signed integer $a$.

Before binomial authentication, \eqref{eq:large-x} and $R=16^t$
already give $x>2\ell(t+1)$ and $x>2t$. For the actual
representatives, $\max(1,\abs{A_j})\le\sqrt H$ and likewise for
$C$. Thus Lemma~\ref{lem:signed-bound} bounds every relevant
coefficient of an ordinary range by $H^{x/2}$. For the derived label,
the same bound follows from $\ell(t+1)<x/2$. For the main binomial,
$R+t<H$ and $t<x/2$ give it again. Consequently, for all eleven
digits and all relevant positions,
\[
 \abs{\binom nj}\le H^{x/2}\le\sqrt q.
\]
As $q$ is already far greater than 4096,
\[
 \frac{\abs{r_0}}{q^k}
 \le\frac{\sqrt q}{q-1}<\frac1{32}.
\]
All common-norm reconstructions are therefore valid, including those
with negative upper indices.

For the main digit, both $B$ and $\binom{R+t}{t}$ are positive
and less than $q/2$: the former follows from $B\le BH$ and
$q=(BH)^x$ with $x\ge3$, and the latter from its square-root bound.
Their congruence modulo $q$ is therefore equality. This authenticates
$B=\binom{R+t}{t}$.

\subsection{Decode rows and recover a finite refutation}

The range congruences have the form
\[
 N\mid\binom Ct,\qquad
 N\mid\binom{A_j}{b_j+1},\qquad
 N\mid\LL(C,\boldsymbol A),
\]
with an additional range for the polynomially derived label. For row
$i-1$, choose $p^e$ from Lemma~\ref{lem:selectors}. By
Lemma~\ref{lem:signed-range}, the ordinary residue of $C$ modulo
$p^{xe}$ lies in $[0,t-1]$. Equation~\eqref{eq:index} identifies
that residue modulo $p^e$ with $i-1$. Since $p^e>t$, it is exactly
$i-1$. All other columns recover their bounded nonnegative local
values in the same way.

The inequalities in \eqref{eq:large-x} give
$2+t+x+y+B<2x$. The elementary inequality
$(2x)^{2^{23}}<2^x$ for $x>2^{32}$ then bounds every decoded local
value by $2^x\le p^{xe}$. Its divisibility by $p^{xe}$ forces it to
vanish. For the derived label, the decoded syntax inputs lie in
$[0,t]$. Apply the same structural ledger with input bound $t+1$ to
obtain
\[
 0\le J(\text{local syntax})\le(t+1)^\ell
 <(2+t+x+y+B)^E<2^x.
\]
Its small range residue and its locally evaluated value are congruent
modulo $p^{xe}$ and both lie in $[0,p^{xe})$, so they agree.
Every row is valid.

The recovered beta accumulator starts at $B+1$ and ends at 1.
The inequality $y>B$ makes the initial remainder proper without a
boundary exception. Valid local updates have positive weights $1+jB$;
starting from $B+1>0$, all accumulators are positive. Multiply the
updates and cancel them to obtain \eqref{eq:balance}. Lemma~\ref{lem:binomial-growth}
and the flow argument then give a finite $C_{222}$ refutation. The
retained logical calibration gives a finite ZFC refutation.

\section{Constructive completeness}

Start with a finite $C_{222}$ refutation, expand it to a proof tree,
and choose $t=16+h^2$ beyond every required row count and label bound.
Truth padding supplies exactly $t$ rows. Set $R=16^t$ and
$B=\binom{R+t}{t}$. Build the pending accumulators using weights
$1+jB$.

Let $O$ be the odd part of $t!$. Choose an arbitrarily large positive
odd multiple $w$ of $O$ and set $y=RBw^2$. By
Lemma~\ref{lem:binomial-growth}, $B$ is odd; $v_2(R)=4t>v_2(t!)$.
Thus $y$ is an even multiple of $t!$, and can exceed every accumulator.
The beta moduli $1+(i+1)y$, $0\le i\le t$, are pairwise coprime:
a common prime would divide $j-i\le t$ and hence $y$, an impossibility.
They are all odd. CRT supplies the beta code $x$; choose a positive,
sufficiently large even lift. The endpoint quotients and local beta
quotients then satisfy their required bounds. The latter are strictly
below $x$ because their moduli are at least two.

Set $N=B^x$, which is odd. For each full prime power $p^e\Vert B$,
use the compatible row from Lemma~\ref{lem:selectors}, and prescribe
the row's actual values modulo $p^{2xe}$. CRT supplies all nine
representatives modulo $N^2$. The extra precision is essential:
division by any range factorial $k!$, with $k\le x$, loses at most
$k-1<x\le xe$ powers of $p$. The resulting binomial therefore
remains divisible by $p^{xe}$. This proves completeness of the range
conditions without using a false converse to Lemma~\ref{lem:signed-range}.

Choose the dominant schema representative $a$ odd and the other eight
representatives even. The inherited label constructors give
$J\equiv1+a+b\pmod2$ and $J\ge2a^2$, so the semantic anchor is
even. These parity choices make $H$ odd. If $3\mid B$, then $BH$
is already a multiple of 3. If $3\nmid B$, the modulus $N^2$ is
coprime to 6, and the representatives may additionally be chosen so
that $H\equiv0\pmod3$. Two independent square residues suffice:
their possible sums modulo 3 are 0, 1, and 2. This adjustment preserves
the stated parities and the CRT classes.

Choose the other representatives positive and at most $6N^2$; take
$a$ in a compatible class between $2N^2$ and $8N^2$. Then
$J\ge8N^4$ dominates twice every other requested index. The front
indices satisfy $R+t<B<N^2$ and $t,y+1,x<N^2$. Increasing the
representatives changes the self-sized radix using their actual values,
so it does not invalidate the coefficient bounds.

Now $q=(BH)^x$ is a positive odd multiple of 3, hence it has the
required affine form $6z+3$. Compute the exact powers in
\eqref{eq:six}, the successor denominator, the actual index atoms, and
their magnitude frame. Choose the lower frame abscissa $L$ and a sufficiently large
$A=L^2+u^2+v^2$ by Theorem~\ref{thm:two-square}, so that
\eqref{eq:common-root} holds. The inherited anchor and pooled
orientation completion at this $A$ supplies all requested Pell pairs.
Use their polynomial values \eqref{eq:mu} as the eleven modular
numerators.

The genuine normalized remainders have absolute value below $M/32$.
Therefore
\[
 4\sum_{i=1}^{11}\rho_i^2<\frac{44}{1024}M^2<M^2.
\]
The remaining nonnegative integer gap in \eqref{eq:norm} has a
$\G$ representation. All quotient witnesses are integral by the
common-root and range congruences. This constructs a zero of the
new raw polynomial system from every finite refutation, without
making any choice depend circularly on the subsequently chosen Pell
center.

\section{The literal polynomial and finite verification}

Every source definition is a quadratic integer equation. The builder
records the exact changes from the pinned parent, the recomputed
magnitude frame, the full power-request census, and every deleted
private coordinate. It then applies the retained Horn factorization
and polynomial contractions.

A permitted monic contraction removes a row of the form
$\epsilon z_i+f=0$, where $\epsilon\in\{1,-1\}$ and $f$ does
not involve $z_i$. Substitution gives an integer polynomial inverse,
and the checker requires every remaining residual to stay quadratic.
Recorded affine consequences and integer residual-row operations
preserve the common zero set. For example, replacing $R_i$ by
$R_i+cR_j$ is reversed by subtraction. Integer coordinate translations
are bijections on the declared domain. The sum-of-squares polynomial
may change under these operations, but its zero set remains equivalent.

% BEGIN report/verification.tex
\begin{center}\small
\begin{tabular}{@{}lr@{}}
\toprule Finite audit & Count \\ \midrule
Literal exact-power requests retained & 6 \\
Modular numerator definitions checked & 11 \\
Principal signed Pell pairs checked & 14 \\
Retained frame atoms & 21 \\
Independent Horn factoring edits & 18 \\
Integral monic restorations & 208 \\
Coordinate translations before final refinements & 69 \\
Final balanced / compact affine projections & 2 / 3 \\
Deliberately corrupted algebra artifacts rejected & 7 \\
\bottomrule\end{tabular}\end{center}
Each collected coefficient is checked independently by summing the literal residual squares.
The normal and optimized-Python checks give the same algebraic results.
The complete packet also reproduces all 41 newly built stage and endpoint JSON files byte for byte.
The two-square-center regressions cover 2,210 completions and 3,128 prime-power cases,
including the counterexample to keeping an arbitrary frame fixed.
% END report/verification.tex

The verification code uses explicit exceptions for failed conditions,
so optimized Python does not disable its checks. The arithmetic
regressions cover signed generalized binomial coefficients, signed
range implications, fixed-base binomial parity and growth, prime-power
row selectors, and common-root modular digits. These finite cases
are useful error detectors; the proofs in the preceding sections are
what establish the unbounded claims.

\subsection{Reproduction and integration}

From a fresh extraction, with Python 3.10 or later:
\begin{verbatim}
python3 -B verify.py
python3 -O -B verify.py
\end{verbatim}
The verifier is read-only unless a receipt path is explicitly supplied.
The archive includes deterministic builders and exact source-to-output
certificates; no network access or external symbolic algebra system is
needed for verification. To rebuild the article with ordinary TeX Live,
run \nolinkurl{build_report.sh}. The standalone TeX source contains
all generated report inputs. PDF metadata may vary by TeX installation;
the exact polynomial data do not.

The contribution can be placed as a new directory in
\nolinkurl{Computability/DiophantineConsistency/con-zfc-diophantine-forms}.
This work does not modify an upstream branch. Source licenses and
attribution notices are retained. The integrity manifest is a local
checksum inventory, not an external signature or a proof certificate.

\section{What remains open}

The delivered result is a complete, fixed polynomial with an ordinary
consistency-equivalence proof. It does not establish the smallest
possible variable count, the smallest coefficient height, or a global
optimal arithmetic circuit. Nor does it make a nine-natural-variable
or ten-integer-variable universality result obsolete. In particular,
Sun's September 2026 ten-integer-variable theorem \cite{sun10}
and the contemporary universal-bound work \cite{bayer} concern
different complexity tradeoffs.

The central reusable result here is the source architecture:
one binomial supplies row selectors and multiplicity separation;
one existing exponent request supplies the needed sign; and one
cofinal center supports all modular numerators. The finite code measures
the complete composition. Further improvements should target the
remaining six exact powers, the private anchor witnesses, or a
different local proof calculus. Merely shortening a displayed
abbreviation, omitting a signed case, or discarding a completeness
condition would not be an improvement.

\begin{thebibliography}{9}
\bibitem{repo} V. Reshetnikov, \emph{ProveIt}, inspected commit
\texttt{b7e35224c04fe1031c0a53f80a395e9b32205022}.
\url{https://github.com/VladimirReshetnikov/ProveIt}.
The exact inherited source and archive hashes are in the supplied provenance ledger.
\bibitem{factorialfree} \emph{Removing factorial authenticity from
Diophantine consistency}, incoming ProveIt research packet,
25 October 2026. Retained in \nolinkurl{upstream/factorial_free/}.
The binomial selector and large-evaluation lemmas are inherited from
this packet; its earlier modular numerator sources are separately credited.
\bibitem{farhi} B. Farhi, \emph{An identity involving the least common
multiple of binomial coefficients and its application},
American Mathematical Monthly 116 (2009), 836--839.
\url{https://arxiv.org/abs/0906.2295}.
\bibitem{sun10} Z.-W. Sun, \emph{Ten unknowns for Hilbert's tenth
problem over the integers}, arXiv:2609.01594v2, 2 September 2026.
\url{https://arxiv.org/abs/2609.01594}.
\bibitem{bayer} J. Bayer, M. David, M. Ha\ss ler, Y. Matiyasevich,
and D. Schleicher, \emph{Diophantine equations over $\Z$: universal
bounds and parallel formalization}, arXiv:2506.20909v3,
22 October 2026. \url{https://arxiv.org/html/2506.20909v3}.
\end{thebibliography}

\clearpage\appendix
\section{Complete principal residual list}\label{app:residuals}

The only unknowns are $z_0,\ldots,z_{\MainLastVar}\in\Z$.
Define $P=\sum_{i=0}^{\MainLastRow}R_i^2$ using the following
literal list. There are no additional quantified symbols. This list
and the declared integer domain specify the entire equation; the
JSON file uses the same zero-based coordinate ordering.

% BEGIN report/residuals.tex
\begingroup
\small
\setlength{\jot}{3pt}
\allowdisplaybreaks[4]
\begin{align*}
R_{0}={}& z_{3} - z_{4} - 2z_{4}\,z_{228} \\[2pt]
R_{1}={}& z_{0} - z_{0}\,z_{2} - z_{2} + z_{7} \\[2pt]
R_{2}={}& -1 - z_{2}\,z_{5} + z_{6}\,z_{7} - 2z_{228} \\[2pt]
R_{3}={}& -1 + z_{1} - z_{6}\,z_{7} \\[2pt]
R_{4}={}& 1 + z_{4} + z_{16} - z_{18} - 2z_{18}\,z_{228} \\[2pt]
R_{5}={}& -1 - z_{8}^{2} - 2z_{8}\,z_{9} - z_{9}^{2} + z_{19} \\[2pt]
R_{6}={}& -10z_{8} + z_{8}^{2} - 6z_{8}\,z_{9} - 8z_{8}\,z_{19} - 20z_{9} - 11z_{9}^{2} - 16z_{9}\,z_{19} \\*
&\quad  - 15z_{19} - 4z_{19}^{2} + z_{20} \\[2pt]
R_{7}={}& -2 + 2z_{8} - z_{8}^{2} - 4z_{8}\,z_{20} - 4z_{20} - 4z_{20}^{2} + z_{21} \\[2pt]
R_{8}={}& 1 + 6z_{8} - 36z_{8}^{2} - 24z_{8}\,z_{21} + 2z_{21} - 4z_{21}^{2} + z_{22} \\[2pt]
R_{9}={}& -2 - z_{9} - z_{9}^{2} - 2z_{9}\,z_{10} - z_{10} - z_{10}^{2} + z_{23} \\[2pt]
R_{10}={}& -8 + 7z_{8} - z_{8}^{2} + 4z_{8}\,z_{10} - 4z_{8}\,z_{23} - 4z_{10} + 4z_{23} + z_{24} - z_{46} \\[2pt]
R_{11}={}& -1 + z_{8} - z_{8}^{2} - 2z_{8}\,z_{10} + z_{10} - z_{10}^{2} + z_{25} \\[2pt]
R_{12}={}& 23 + 16z_{8} - 8z_{8}\,z_{9} - 8z_{8}\,z_{10} - 16z_{8}\,z_{19} - 8z_{8}\,z_{25} + 19z_{9} \\*
&\quad  + 9z_{9}^{2} - 36z_{9}\,z_{10} - 12z_{9}\,z_{25} + 16z_{10} + 8z_{10}^{2} - 24z_{10}\,z_{19} - 8z_{19}\,z_{25} \\*
&\quad  - z_{20} + z_{26} - z_{31} \\[2pt]
R_{13}={}& -4z_{24}^{2} - 8z_{24}\,z_{26} - 4z_{26}^{2} + z_{27} \\[2pt]
R_{14}={}& -2 - 8z_{8} + 8z_{8}\,z_{23} + 8z_{8}\,z_{25} - 8z_{9} + 8z_{9}\,z_{23} + 8z_{9}\,z_{25} - 4z_{19} \\*
&\quad  + 8z_{23} - 4z_{23}^{2} - 8z_{23}\,z_{25} + 8z_{25} - 4z_{25}^{2} + z_{28} \\[2pt]
R_{15}={}& 6 + 4z_{8} + 32z_{8}\,z_{9} + 16z_{8}\,z_{19} + 16z_{8}\,z_{25} + 8z_{8}\,z_{28} + 24z_{9} \\*
&\quad  - 32z_{9}\,z_{25} - 16z_{9}\,z_{28} + 12z_{19} - 16z_{19}\,z_{25} - 8z_{19}\,z_{28} - z_{20} + 8z_{25} \\*
&\quad  - 16z_{25}^{2} - 16z_{25}\,z_{28} + 4z_{28} - 4z_{28}^{2} + z_{29} \\[2pt]
R_{16}={}& 22 - 5z_{8} - z_{8}^{2} - 8z_{8}\,z_{9} - 12z_{8}\,z_{10} - 4z_{8}\,z_{23} - 8z_{9} - 16z_{9}\,z_{10} \\*
&\quad  - 16z_{9}\,z_{23} - 20z_{10} - 20z_{10}^{2} - 24z_{10}\,z_{23} - 28z_{23} - 4z_{23}^{2} + z_{30} \\[2pt]
R_{17}={}& 19 - 25z_{8} - 7z_{8}^{2} - 8z_{8}\,z_{9} - 30z_{8}\,z_{10} - 16z_{8}\,z_{25} + 3z_{9} + 5z_{9}^{2} \\*
&\quad  - 4z_{9}\,z_{25} - 27z_{10} - 21z_{10}^{2} - 24z_{10}\,z_{25} - 6z_{23} - 17z_{25} - 4z_{25}^{2} + z_{31} \\[2pt]
R_{18}={}& 2 - 4z_{30}^{2} - 8z_{30}\,z_{31} - 4z_{31}^{2} + z_{32} \\[2pt]
R_{19}={}& -19z_{8} + 8z_{8}^{2} + 8z_{8}\,z_{19} + 8z_{8}\,z_{23} - 46z_{9} - 24z_{9}^{2} + 24z_{9}\,z_{19} \\*
&\quad  + 24z_{9}\,z_{23} - 2z_{10} + 4z_{19} - 4z_{19}^{2} - 8z_{19}\,z_{23} + 14z_{23} - 4z_{23}^{2} - z_{25} \\*
&\quad  + z_{33} \\[2pt]
R_{20}={}& -3 - 6z_{8} - 4z_{10} - 2z_{25} - 4z_{33} - 4z_{33}^{2} + z_{34} \\[2pt]
R_{21}={}& 8 + 42z_{8} + 7z_{8}^{2} + 28z_{8}\,z_{20} + 48z_{9} - 32z_{9}^{2} - 32z_{9}\,z_{19} \\*
&\quad  + 64z_{9}\,z_{20} - 40z_{19}^{2} + 32z_{19}\,z_{20} + 14z_{20} - 4z_{20}^{2} + z_{35} + 2z_{38} \\[2pt]
R_{22}={}& 1 + 6z_{8} - 36z_{8}^{2} - 24z_{8}\,z_{35} + 2z_{35} - 4z_{35}^{2} + z_{36} \\[2pt]
R_{23}={}& -z_{9}^{2} - 4z_{9}\,z_{36} - 4z_{36}^{2} + z_{37} \\[2pt]
R_{24}={}& -6 - 2z_{8} - 4z_{9} - 14z_{19} - 16z_{19}^{2} + z_{38} \\[2pt]
R_{25}={}& 24 + 21z_{9} - 21z_{9}^{2} - 18z_{9}\,z_{10} - 24z_{9}\,z_{23} - 24z_{9}\,z_{38} + 3z_{10} \\*
&\quad  - 9z_{10}^{2} - 16z_{10}\,z_{38} + 3z_{23} - 12z_{23}^{2} - 8z_{23}\,z_{38} + 4z_{38} - 4z_{38}^{2} + z_{39} \\*
&\quad  + 2z_{46} \\[2pt]
R_{26}={}& 16 - 10z_{8} - 20z_{8}^{2} - 16z_{8}\,z_{10} - 24z_{8}\,z_{25} - 24z_{8}\,z_{39} - 12z_{10} \\*
&\quad  - 16z_{10}\,z_{25} - 16z_{10}\,z_{39} - 14z_{25} - 4z_{25}^{2} - 8z_{25}\,z_{39} + 4z_{39} - 4z_{39}^{2} \\*
&\quad  + z_{40} \\[2pt]
R_{27}={}& -z_{10}^{2} - 2z_{10}\,z_{11} - z_{11}^{2} + z_{41} \\[2pt]
R_{28}={}& 2 + z_{8}^{2} - 4z_{8}\,z_{41} + 2z_{10} - 2z_{10}^{2} - 8z_{10}\,z_{41} - 2z_{25} + 4z_{41} - 4z_{41}^{2} \\*
&\quad  + z_{42} \\[2pt]
R_{29}={}& -31 - 16z_{10} - 16z_{10}\,z_{23} + 16z_{10}\,z_{25} - 8z_{23} + 12z_{23}^{2} - 8z_{23}\,z_{25} + 24z_{25} \\*
&\quad  - 4z_{25}^{2} + z_{43} - 4z_{46} \\[2pt]
R_{30}={}& -4z_{10} - 16z_{25}^{2} - 32z_{25}\,z_{41} - 16z_{25}\,z_{42} - 16z_{25}\,z_{43} - 4z_{41} - 16z_{41}^{2} \\*
&\quad  - 16z_{41}\,z_{42} - 16z_{41}\,z_{43} - 2z_{42} - 4z_{42}^{2} - 8z_{42}\,z_{43} - 4z_{43}^{2} + z_{44} \\[2pt]
R_{31}={}& -16z_{8}^{2} + 8z_{8}\,z_{10} + 32z_{8}\,z_{25} + 16z_{8}\,z_{28} + z_{10} - z_{10}^{2} - 8z_{10}\,z_{25} \\*
&\quad  - 4z_{10}\,z_{28} - 16z_{25}^{2} - 16z_{25}\,z_{28} - 4z_{28}^{2} + z_{45} \\[2pt]
R_{32}={}& -1 - 8z_{10} - 4z_{10}^{2} + 8z_{10}\,z_{23} + 8z_{23} - 4z_{23}^{2} + z_{46} \\[2pt]
R_{33}={}& -1 - 9z_{10}^{2} + 24z_{10}\,z_{23} + 12z_{10}\,z_{46} - 16z_{23}^{2} - 16z_{23}\,z_{46} - 4z_{46}^{2} + z_{47} \\[2pt]
R_{34}={}& 1 + z_{9} - z_{9}^{2} - 4z_{9}\,z_{10} - 4z_{9}\,z_{47} - 4z_{10}^{2} - 8z_{10}\,z_{47} - 4z_{47}^{2} + z_{48} \\[2pt]
R_{35}={}& -1 - 4z_{48} - 4z_{48}^{2} + z_{49} \\[2pt]
R_{36}={}& -1 - z_{10}^{2} - 8z_{10}\,z_{48} - 4z_{10}\,z_{49} - 16z_{48}^{2} - 16z_{48}\,z_{49} - 4z_{49}^{2} + z_{50} \\[2pt]
R_{37}={}& 2 + z_{9} - z_{9}^{2} - 4z_{9}\,z_{10} - 4z_{9}\,z_{50} - 4z_{10}^{2} - 8z_{10}\,z_{50} - 4z_{50}^{2} + z_{51} \\[2pt]
R_{38}={}& -4z_{45}^{2} - 8z_{45}\,z_{51} - 4z_{51}^{2} + z_{52} \\[2pt]
R_{39}={}& -9 + 7z_{8} - z_{8}^{2} - 8z_{8}\,z_{45} - 4z_{8}\,z_{52} + 24z_{45} - 16z_{45}^{2} - 16z_{45}\,z_{52} \\*
&\quad  + 12z_{52} - 4z_{52}^{2} + z_{53} \\[2pt]
R_{40}={}& -1 - 4z_{53} - 4z_{53}^{2} + z_{54} \\[2pt]
R_{41}={}& -z_{10}^{2} - 8z_{10}\,z_{53} - 4z_{10}\,z_{54} - 16z_{53}^{2} - 16z_{53}\,z_{54} - 4z_{54}^{2} + z_{55} \\[2pt]
R_{42}={}& -z_{8} - z_{8}^{2} - 4z_{8}\,z_{10} - 4z_{8}\,z_{55} - 4z_{10} - 4z_{10}^{2} - 8z_{10}\,z_{55} - 4z_{55} \\*
&\quad  - 4z_{55}^{2} + z_{56} \\[2pt]
R_{43}={}& -10 - 12z_{56} - 4z_{56}^{2} + z_{57} \\[2pt]
R_{44}={}& 2 + 4z_{10} + 16z_{10}\,z_{56} + 8z_{10}\,z_{57} - 4z_{23} - 16z_{23}\,z_{56} - 8z_{23}\,z_{57} - z_{46} \\*
&\quad  + 8z_{56} - 16z_{56}^{2} - 16z_{56}\,z_{57} + 4z_{57} - 4z_{57}^{2} + z_{58} \\[2pt]
R_{45}={}& z_{9} - z_{9}^{2} + 8z_{9}\,z_{10} - 8z_{9}\,z_{23} - 4z_{9}\,z_{58} - 16z_{10}^{2} + 32z_{10}\,z_{23} \\*
&\quad  + 16z_{10}\,z_{58} - 16z_{23}^{2} - 16z_{23}\,z_{58} - 4z_{58}^{2} + z_{59} \\[2pt]
R_{46}={}& -10 - 12z_{51} - 4z_{51}^{2} + z_{60} \\[2pt]
R_{47}={}& 1 - 16z_{51}^{2} - 16z_{51}\,z_{59} - 16z_{51}\,z_{60} - 4z_{59}^{2} - 8z_{59}\,z_{60} - 4z_{60}^{2} + z_{61} \\[2pt]
R_{48}={}& z_{10} - z_{10}^{2} - 8z_{10}\,z_{59} - 4z_{10}\,z_{61} - 16z_{59}^{2} - 16z_{59}\,z_{61} - 4z_{61}^{2} + z_{62} \\[2pt]
R_{49}={}& -6z_{8} - 3z_{8}^{2} - 18z_{8}\,z_{9} - 16z_{9} - 27z_{9}^{2} + 4z_{9}\,z_{10} + 3z_{10} + 3z_{10}^{2} \\*
&\quad  - 4z_{10}\,z_{19} - 19z_{19} - 8z_{19}^{2} + z_{20} + z_{63} - 4z_{76} \\[2pt]
R_{50}={}& 1 - 8z_{8} - 4z_{8}\,z_{10} + 8z_{8}\,z_{63} - 16z_{9} - 8z_{9}\,z_{10} + 16z_{9}\,z_{63} + z_{10} \\*
&\quad  + z_{10}^{2} + 4z_{10}\,z_{19} + 8z_{19} - 8z_{19}\,z_{63} + 3z_{63} - 4z_{63}^{2} + z_{64} \\[2pt]
R_{51}={}& -2 - 12z_{8} - 16z_{8}\,z_{10} + 16z_{8}\,z_{62} + 16z_{8}\,z_{64} - 24z_{9} - 32z_{9}\,z_{10} \\*
&\quad  + 32z_{9}\,z_{62} + 32z_{9}\,z_{64} + 4z_{10} + 4z_{10}^{2} + 16z_{10}\,z_{19} + 12z_{19} - 16z_{19}\,z_{62} \\*
&\quad  - 16z_{19}\,z_{64} - 4z_{62}^{2} - 8z_{62}\,z_{64} - 4z_{63} - 2z_{64} - 4z_{64}^{2} + z_{65} \\[2pt]
R_{52}={}& -2 - 4z_{65} - 4z_{65}^{2} + z_{66} \\[2pt]
R_{53}={}& -z_{13} - 2z_{13}\,z_{228} + z_{67} \\[2pt]
R_{54}={}& z_{12} - z_{13} - 2z_{22}\,z_{67} - z_{68} \\[2pt]
R_{55}={}& z_{12} - z_{13} - 4z_{24}\,z_{67} - 2z_{27}\,z_{67} - z_{69} \\[2pt]
R_{56}={}& -2z_{29}\,z_{67} + 2z_{44}\,z_{67} - 2z_{67} - z_{70} + z_{75} \\[2pt]
R_{57}={}& z_{12} - z_{13} - 4z_{30}\,z_{67} - 2z_{32}\,z_{67} - z_{71} \\[2pt]
R_{58}={}& z_{12} - z_{13} - 2z_{34}\,z_{67} + 6z_{67} - z_{72} \\[2pt]
R_{59}={}& z_{12} - z_{13} - 4z_{36}\,z_{67} - 2z_{37}\,z_{67} - z_{73} \\[2pt]
R_{60}={}& z_{12} - z_{13} - 2z_{40}\,z_{67} - z_{74} \\[2pt]
R_{61}={}& z_{12} - z_{13} - 2z_{44}\,z_{67} - 4z_{67} - z_{75} \\[2pt]
R_{62}={}& -z_{8}\,z_{9} + z_{8}\,z_{10} + z_{9} + z_{9}\,z_{10} - z_{10} - z_{10}^{2} + z_{76} \\[2pt]
R_{63}={}& 1 + 2z_{11} + z_{11}^{2} - z_{76}^{2} + z_{77} \\[2pt]
R_{64}={}& -z_{77}^{2} + z_{78} \\[2pt]
R_{65}={}& -4z_{65}\,z_{67} - 2z_{66}\,z_{67} + 2z_{67} + z_{79} \\[2pt]
R_{66}={}& -z_{12}^{2} + 2z_{12}\,z_{13} + 2z_{12}\,z_{79} - z_{13}^{2} - 2z_{13}\,z_{79} - z_{79}^{2} + z_{80} \\[2pt]
R_{67}={}& -1 - z_{9} - 2z_{9}\,z_{228} + z_{81} \\[2pt]
R_{68}={}& -z_{13}\,z_{81} + z_{83}\,z_{84} + z_{85} \\[2pt]
R_{69}={}& -z_{8} - 2z_{8}\,z_{228} + z_{82} + 2z_{228} \\[2pt]
R_{70}={}& 1 + 2z_{8} + 4z_{8}\,z_{228} + 4z_{9} + 8z_{9}\,z_{228} - 2z_{19} - 4z_{19}\,z_{228} + z_{83} + 4z_{228} \\[2pt]
R_{71}={}& -z_{12}\,z_{82} + z_{84} \\[2pt]
R_{72}={}& z_{12}\,z_{81} + 6z_{13} - 8z_{13}\,z_{81} - 4z_{13}\,z_{82} - z_{13}\,z_{83} - 12z_{13}\,z_{228} + z_{86} \\[2pt]
R_{73}={}& -z_{12}\,z_{68} + z_{13}\,z_{68} + z_{87} \\[2pt]
R_{74}={}& -z_{69}\,z_{87} + z_{88} \\[2pt]
R_{75}={}& -z_{70}\,z_{88} + z_{89} \\[2pt]
R_{76}={}& -z_{71}\,z_{89} + z_{90} \\[2pt]
R_{77}={}& -z_{72}\,z_{90} + z_{91} \\[2pt]
R_{78}={}& -z_{73}\,z_{91} + z_{92} \\[2pt]
R_{79}={}& -z_{74}\,z_{92} + z_{93} \\[2pt]
R_{80}={}& -z_{75}\,z_{93} + z_{94} \\[2pt]
R_{81}={}& -z_{85}\,z_{94} + z_{95} \\[2pt]
R_{82}={}& -z_{86}\,z_{95} + z_{96} \\[2pt]
R_{83}={}& -z_{96}^{2} + z_{97} \\[2pt]
R_{84}={}& -z_{2} - z_{2}\,z_{16} + z_{16} + z_{98} \\[2pt]
R_{85}={}& -z_{14}\,z_{98} + z_{99} \\[2pt]
R_{86}={}& -2z_{1}^{2} + 2z_{1}\,z_{12} + 2z_{1}\,z_{13} + 2z_{1}\,z_{99} + 2z_{1}\,z_{100} - z_{12}^{2} - 2z_{12}\,z_{99} \\*
&\quad  - z_{13}^{2} - 2z_{13}\,z_{100} + z_{17}\,z_{102} - z_{78}\,z_{97} - z_{80}\,z_{97} - z_{99}^{2} - z_{100}^{2} \\[2pt]
R_{87}={}& -z_{2}\,z_{15} + z_{15} - z_{15}\,z_{98} + z_{100} \\[2pt]
R_{88}={}& -z_{8} - z_{9} - 4z_{19} - 2z_{22} - 4z_{24} - 2z_{27} - 2z_{29} - 4z_{30} - 2z_{32} - 2z_{34} \\*
&\quad  - 4z_{36} - 2z_{37} - 2z_{40} - 2z_{44} - 4z_{65} - 2z_{66} + z_{101} \\[2pt]
R_{89}={}& -3z_{105} - 6z_{105}\,z_{221} + z_{107} \\[2pt]
R_{90}={}& z_{103}\,z_{212} - z_{106} - z_{107}\,z_{128} - z_{128} - 2z_{128}\,z_{228} \\[2pt]
R_{91}={}& z_{108}\,z_{213} - z_{110}\,z_{131} - z_{111} \\[2pt]
R_{92}={}& z_{112}\,z_{213} - z_{113}\,z_{131} - z_{114} \\[2pt]
R_{93}={}& z_{115}\,z_{213} - z_{116}\,z_{131} - z_{117} \\[2pt]
R_{94}={}& z_{118}\,z_{213} - z_{119}\,z_{131} - z_{120} \\[2pt]
R_{95}={}& z_{121}\,z_{214} - z_{123}\,z_{131} - z_{124} \\[2pt]
R_{96}={}& z_{125}\,z_{214} - z_{126}\,z_{131} - z_{127} \\[2pt]
R_{97}={}& -z_{17}\,z_{128} + z_{131} \\[2pt]
R_{98}={}& -z_{131}\,z_{135} + z_{134}\,z_{212} - z_{136} \\[2pt]
R_{99}={}& -z_{131}\,z_{138} + z_{137}\,z_{213} - z_{139} \\[2pt]
R_{100}={}& z_{140} - 2z_{140}\,z_{142} + z_{143} \\[2pt]
R_{101}={}& -z_{140}\,z_{143} - z_{143} + z_{144} \\[2pt]
R_{102}={}& -1 + z_{143}^{2} - z_{144}^{2} + z_{224}^{2} \\[2pt]
R_{103}={}& -z_{141}\,z_{146} - 7z_{146} + z_{147} \\[2pt]
R_{104}={}& 2z_{145} + z_{145}^{2} + z_{146}^{2} - z_{147}^{2} \\[2pt]
R_{105}={}& 1 - z_{4} - 32z_{141}\,z_{148} + z_{145} + 16z_{146} - z_{147} + 33z_{148} \\[2pt]
R_{106}={}& -12 - 2z_{141} - 4z_{141}\,z_{228} + z_{149} - 24z_{228} + 4z_{228}^{2} \\[2pt]
R_{107}={}& z_{17} + z_{149}\,z_{150} - z_{183} - z_{184} - 2z_{184}\,z_{228} + z_{185} \\[2pt]
R_{108}={}& -1 + 2z_{141}\,z_{196} - z_{151} + 14z_{196} - z_{196}^{2} \\[2pt]
R_{109}={}& 3 + z_{151}\,z_{152} - z_{183} - z_{184}\,z_{196} + z_{185} + 6z_{221} \\[2pt]
R_{110}={}& -z_{141}\,z_{154} - 7z_{154} + z_{155} \\[2pt]
R_{111}={}& 2z_{153} + z_{153}^{2} + z_{154}^{2} - z_{155}^{2} \\[2pt]
R_{112}={}& -39 + z_{131}\,z_{229} - 8z_{141} - 12z_{141}\,z_{221} - 36z_{221} + 36z_{221}^{2} \\[2pt]
R_{113}={}& 1 - z_{103} - z_{141}\,z_{154} + z_{153} - 3z_{154} + 6z_{154}\,z_{221} \\[2pt]
R_{114}={}& 7 - z_{131}\,z_{229} + 2z_{141} + z_{156} - 12z_{221} \\[2pt]
R_{115}={}& -1 + z_{104} - z_{145} - 3z_{146} - 6z_{146}\,z_{221} + z_{147} + z_{156}\,z_{157} \\[2pt]
R_{116}={}& 7 + z_{141} - z_{141}\,z_{159} - 7z_{159} + z_{160} \\[2pt]
R_{117}={}& 1 + 2z_{158} + z_{158}^{2} - 2z_{159} + z_{159}^{2} - z_{160}^{2} \\[2pt]
R_{118}={}& 4 - z_{108} + z_{141} - z_{141}\,z_{159} + z_{158} - 3z_{159} + 6z_{159}\,z_{221} - 6z_{221} \\[2pt]
R_{119}={}& -3z_{104} - 6z_{104}\,z_{221} + z_{109} \\[2pt]
R_{120}={}& -z_{141}\,z_{162} - 7z_{162} + z_{163} \\[2pt]
R_{121}={}& 2z_{161} + z_{161}^{2} + z_{162}^{2} - z_{163}^{2} \\[2pt]
R_{122}={}& 1 - z_{112} - z_{141}\,z_{162} + z_{161} - 3z_{162} + 6z_{162}\,z_{221} \\[2pt]
R_{123}={}& -z_{141}\,z_{165} - 7z_{165} + z_{166} \\[2pt]
R_{124}={}& 2z_{164} + z_{164}^{2} + z_{165}^{2} - z_{166}^{2} \\[2pt]
R_{125}={}& 1 - z_{115} - z_{141}\,z_{165} + z_{164} - 3z_{165} + 6z_{165}\,z_{221} \\[2pt]
R_{126}={}& -z_{141}\,z_{168} - 7z_{168} + z_{169} \\[2pt]
R_{127}={}& 2z_{167} + z_{167}^{2} + z_{168}^{2} - z_{169}^{2} \\[2pt]
R_{128}={}& 1 - z_{118} - z_{141}\,z_{168} + z_{167} - 3z_{168} + 6z_{168}\,z_{221} \\[2pt]
R_{129}={}& -z_{141}\,z_{171} - 7z_{171} + z_{172} \\[2pt]
R_{130}={}& 2z_{170} + z_{170}^{2} + z_{171}^{2} - z_{172}^{2} \\[2pt]
R_{131}={}& 1 - z_{121} - z_{141}\,z_{171} + z_{170} - 3z_{171} + 6z_{171}\,z_{221} \\[2pt]
R_{132}={}& -z_{141}\,z_{174} - 7z_{174} + z_{175} \\[2pt]
R_{133}={}& 2z_{173} + z_{173}^{2} + z_{174}^{2} - z_{175}^{2} \\[2pt]
R_{134}={}& -1 + z_{122} + z_{156}\,z_{176} - z_{173} - 3z_{174} - 6z_{174}\,z_{221} + z_{175} \\[2pt]
R_{135}={}& -z_{141}\,z_{178} - 7z_{178} + z_{179} \\[2pt]
R_{136}={}& 2z_{177} + z_{177}^{2} + z_{178}^{2} - z_{179}^{2} \\[2pt]
R_{137}={}& 1 - z_{125} - z_{141}\,z_{178} + z_{177} - 3z_{178} + 6z_{178}\,z_{221} \\[2pt]
R_{138}={}& -z_{141}\,z_{181} - 7z_{181} + z_{182} \\[2pt]
R_{139}={}& 2z_{180} + z_{180}^{2} + z_{181}^{2} - z_{182}^{2} \\[2pt]
R_{140}={}& 1 - z_{129}\,z_{131} - z_{130} - z_{141}\,z_{181} + z_{180} - 3z_{181} + 6z_{181}\,z_{221} \\[2pt]
R_{141}={}& 1 - z_{141}\,z_{184} - 7z_{184} + z_{185} \\[2pt]
R_{142}={}& -1 + 2z_{183} + z_{183}^{2} + z_{184}^{2} - 2z_{185} - z_{185}^{2} \\[2pt]
R_{143}={}& z_{128} + z_{156}\,z_{186} - z_{183} - 3z_{184} - 6z_{184}\,z_{221} + z_{185} \\[2pt]
R_{144}={}& -z_{141}\,z_{188} - 7z_{188} + z_{189} \\[2pt]
R_{145}={}& 2z_{187} + z_{187}^{2} + z_{188}^{2} - z_{189}^{2} \\[2pt]
R_{146}={}& 1 - z_{131}\,z_{132} - z_{133} - z_{141}\,z_{188} + z_{187} - 3z_{188} + 6z_{188}\,z_{221} \\[2pt]
R_{147}={}& -z_{141}\,z_{191} - 7z_{191} + z_{192} \\[2pt]
R_{148}={}& 2z_{190} + z_{190}^{2} + z_{191}^{2} - z_{192}^{2} \\[2pt]
R_{149}={}& 1 - z_{134} - z_{141}\,z_{191} + z_{190} - 3z_{191} + 6z_{191}\,z_{221} \\[2pt]
R_{150}={}& -z_{141}\,z_{194} - 7z_{194} + z_{195} \\[2pt]
R_{151}={}& 2z_{193} + z_{193}^{2} + z_{194}^{2} - z_{195}^{2} \\[2pt]
R_{152}={}& 1 - z_{137} - z_{141}\,z_{194} + z_{193} - 3z_{194} + 6z_{194}\,z_{221} \\[2pt]
R_{153}={}& -12 - 2z_{0}^{2} - 2z_{0}\,z_{4} - z_{1}^{2} - z_{2}^{2} - z_{4}^{2} + 2z_{8} + 2z_{8}\,z_{9} - z_{10}^{2} \\*
&\quad  - z_{11}^{2} - z_{12}^{2} - z_{13}^{2} - z_{14}^{2} - z_{15}^{2} - z_{16}^{2} - z_{17}^{2} - z_{19} - z_{101}^{2} - z_{104}^{2} \\*
&\quad  - z_{122}^{2} - z_{128}^{2} + z_{140} - z_{196}^{2} - 36z_{221} - 36z_{221}^{2} - 4z_{228} - 4z_{228}^{2} \\[2pt]
R_{154}={}& z_{196} - z_{227} - 2z_{227}\,z_{228} \\[2pt]
R_{155}={}& -16 + z_{0} - z_{197}^{2} \\[2pt]
R_{156}={}& -z_{141}\,z_{199} - 7z_{199} + z_{208} \\[2pt]
R_{157}={}& -1 + z_{198}^{2} + z_{199}^{2} - z_{208}^{2} \\[2pt]
R_{158}={}& -z_{201} - z_{201}\,z_{202} + z_{209} \\[2pt]
R_{159}={}& -1 + z_{200}^{2} + z_{201}^{2} - z_{209}^{2} \\[2pt]
R_{160}={}& -z_{194}^{2} + z_{210} \\[2pt]
R_{161}={}& z_{199} - 2z_{204}\,z_{210} - z_{210} \\[2pt]
R_{162}={}& -4z_{199} - 8z_{199}\,z_{203} + z_{202} \\[2pt]
R_{163}={}& -6 - z_{141} - z_{198}\,z_{205} + z_{202} \\[2pt]
R_{164}={}& -1 - z_{193} - z_{198}\,z_{206} + z_{200} \\[2pt]
R_{165}={}& -z_{101} - 4z_{194}\,z_{207} + z_{201} \\[2pt]
R_{166}={}& -2 + 2z_{0}^{2} + 2z_{0}\,z_{4} - 2z_{0}\,z_{146} - 2z_{0}\,z_{154} + z_{1}^{2} - 2z_{1}\,z_{184} + z_{2}^{2} \\*
&\quad  - 2z_{2}\,z_{174} + z_{4}^{2} - 2z_{4}\,z_{154} + z_{8}^{2} - 2z_{8}\,z_{159} + z_{9}^{2} - 2z_{9}\,z_{162} + z_{10}^{2} \\*
&\quad  - 2z_{10}\,z_{165} + z_{11}^{2} - 2z_{11}\,z_{168} + z_{12}^{2} - 2z_{12}\,z_{171} + z_{13}^{2} - 2z_{13}\,z_{178} + z_{14}^{2} \\*
&\quad  - 2z_{14}\,z_{181} + z_{15}^{2} - 2z_{15}\,z_{188} + z_{16}^{2} - 2z_{16}\,z_{191} - 2z_{101}\,z_{194} - z_{141}\,z_{211} \\*
&\quad  + z_{145}^{2} + z_{153}^{2} + z_{158}^{2} + 2z_{159} + z_{161}^{2} + z_{164}^{2} + z_{167}^{2} + z_{170}^{2} + z_{173}^{2} + z_{177}^{2} \\*
&\quad  + z_{180}^{2} + z_{183}^{2} + z_{187}^{2} + z_{190}^{2} + z_{193}^{2} + 2z_{194}^{2} - 6z_{211} - z_{215}^{2} - z_{216}^{2} - z_{217} \\*
&\quad  - z_{217}^{2} \\[2pt]
R_{167}={}& -z_{104}\,z_{212} + z_{128} \\[2pt]
R_{168}={}& -z_{109}\,z_{213} + z_{128} \\[2pt]
R_{169}={}& -z_{122}\,z_{214} + z_{128} \\[2pt]
R_{170}={}& -1 + z_{2} - z_{3}\,z_{222} \\[2pt]
R_{171}={}& -2 - z_{101}^{2} - z_{146}^{2} - z_{154}^{2} + 2z_{159} - z_{159}^{2} - z_{162}^{2} - z_{165}^{2} - z_{168}^{2} - z_{171}^{2} \\*
&\quad  - z_{174}^{2} - z_{178}^{2} - z_{181}^{2} - z_{184}^{2} - z_{188}^{2} - z_{191}^{2} + z_{194}^{2} - z_{215}^{2} - z_{216}^{2} - z_{217} \\*
&\quad  - z_{217}^{2} \\[2pt]
R_{172}={}& -1 - 4z_{106}^{2} - 4z_{111}^{2} - 4z_{114}^{2} - 4z_{117}^{2} - 4z_{120}^{2} - 4z_{124}^{2} - 4z_{127}^{2} + z_{128}^{2} \\*
&\quad  - 4z_{130}^{2} - 4z_{133}^{2} - 4z_{136}^{2} - 4z_{139}^{2} - z_{218}^{2} - z_{219}^{2} - z_{220} - z_{220}^{2} \\[2pt]
R_{173}={}& -1 + z_{222} - 4z_{223} - 4z_{223}^{2} \\[2pt]
R_{174}={}& 7 + z_{141} - z_{224}^{2} - z_{225}^{2} - z_{226}^{2} \\[2pt]
R_{175}={}& -3 - 2z_{0} - 2z_{4} + 2z_{8} - z_{8}^{2} - z_{9}^{2} - z_{10}^{2} - z_{11}^{2} - z_{12}^{2} - z_{13}^{2} \\*
&\quad  - z_{14}^{2} - z_{15}^{2} - z_{16}^{2} + z_{227} \\
\end{align*}
\endgroup
% END report/residuals.tex
\clearpage
\section{Complete compact residual list}\label{app:compact}

For the companion, the only unknowns are
$w_0,\ldots,w_{\CompactLastVar}\in\Z$. Its equation is
\[
 Q(\boldsymbol w)=\sum_{i=0}^{\CompactLastRow}S_i(\boldsymbol w)^2=0.
\]
These are exactly the residuals in
\nolinkurl{equations/compact/equation.json}, with $w$ used as the
printed coordinate name. The compact equation has an integer solution
if and only if a finite syntactic ZFC refutation exists.

% BEGIN report/compact_residuals.tex
\begingroup
\small
\setlength{\jot}{3pt}
\allowdisplaybreaks[4]
\begin{align*}
S_{0}={}& w_{3} - w_{4} - 2w_{4}\,w_{227} \\[2pt]
S_{1}={}& w_{0} - w_{0}\,w_{2} - w_{2} + w_{7} \\[2pt]
S_{2}={}& -1 - w_{2}\,w_{5} + w_{6}\,w_{7} - 2w_{227} \\[2pt]
S_{3}={}& -1 + w_{1} - w_{6}\,w_{7} \\[2pt]
S_{4}={}& 1 + w_{4} + w_{16} - w_{18} - 2w_{18}\,w_{227} \\[2pt]
S_{5}={}& -1 - w_{8}^{2} - 2w_{8}\,w_{9} - w_{9}^{2} + w_{19} \\[2pt]
S_{6}={}& -10w_{8} + w_{8}^{2} - 6w_{8}\,w_{9} - 8w_{8}\,w_{19} - 20w_{9} - 11w_{9}^{2} - 16w_{9}\,w_{19} \\*
&\quad  - 15w_{19} - 4w_{19}^{2} + w_{20} \\[2pt]
S_{7}={}& -2 + 2w_{8} - w_{8}^{2} - 4w_{8}\,w_{20} - 4w_{20} - 4w_{20}^{2} + w_{21} \\[2pt]
S_{8}={}& 1 + 6w_{8} - 36w_{8}^{2} - 24w_{8}\,w_{21} + 2w_{21} - 4w_{21}^{2} + w_{22} \\[2pt]
S_{9}={}& -2 - w_{9} - w_{9}^{2} - 2w_{9}\,w_{10} - w_{10} - w_{10}^{2} + w_{23} \\[2pt]
S_{10}={}& -8 + 7w_{8} - w_{8}^{2} + 4w_{8}\,w_{10} - 4w_{8}\,w_{23} - 4w_{10} + 4w_{23} + w_{24} - w_{46} \\[2pt]
S_{11}={}& -1 + w_{8} - w_{8}^{2} - 2w_{8}\,w_{10} + w_{10} - w_{10}^{2} + w_{25} \\[2pt]
S_{12}={}& 23 + 16w_{8} - 8w_{8}\,w_{9} - 8w_{8}\,w_{10} - 16w_{8}\,w_{19} - 8w_{8}\,w_{25} + 19w_{9} \\*
&\quad  + 9w_{9}^{2} - 36w_{9}\,w_{10} - 12w_{9}\,w_{25} + 16w_{10} + 8w_{10}^{2} - 24w_{10}\,w_{19} - 8w_{19}\,w_{25} \\*
&\quad  - w_{20} + w_{26} - w_{31} \\[2pt]
S_{13}={}& -4w_{24}^{2} - 8w_{24}\,w_{26} - 4w_{26}^{2} + w_{27} \\[2pt]
S_{14}={}& -2 - 8w_{8} + 8w_{8}\,w_{23} + 8w_{8}\,w_{25} - 8w_{9} + 8w_{9}\,w_{23} + 8w_{9}\,w_{25} - 4w_{19} \\*
&\quad  + 8w_{23} - 4w_{23}^{2} - 8w_{23}\,w_{25} + 8w_{25} - 4w_{25}^{2} + w_{28} \\[2pt]
S_{15}={}& 6 + 4w_{8} + 32w_{8}\,w_{9} + 16w_{8}\,w_{19} + 16w_{8}\,w_{25} + 8w_{8}\,w_{28} + 24w_{9} \\*
&\quad  - 32w_{9}\,w_{25} - 16w_{9}\,w_{28} + 12w_{19} - 16w_{19}\,w_{25} - 8w_{19}\,w_{28} - w_{20} + 8w_{25} \\*
&\quad  - 16w_{25}^{2} - 16w_{25}\,w_{28} + 4w_{28} - 4w_{28}^{2} + w_{29} \\[2pt]
S_{16}={}& 22 - 5w_{8} - w_{8}^{2} - 8w_{8}\,w_{9} - 12w_{8}\,w_{10} - 4w_{8}\,w_{23} - 8w_{9} - 16w_{9}\,w_{10} \\*
&\quad  - 16w_{9}\,w_{23} - 20w_{10} - 20w_{10}^{2} - 24w_{10}\,w_{23} - 28w_{23} - 4w_{23}^{2} + w_{30} \\[2pt]
S_{17}={}& 19 - 25w_{8} - 7w_{8}^{2} - 8w_{8}\,w_{9} - 30w_{8}\,w_{10} - 16w_{8}\,w_{25} + 3w_{9} + 5w_{9}^{2} \\*
&\quad  - 4w_{9}\,w_{25} - 27w_{10} - 21w_{10}^{2} - 24w_{10}\,w_{25} - 6w_{23} - 17w_{25} - 4w_{25}^{2} + w_{31} \\[2pt]
S_{18}={}& 2 - 4w_{30}^{2} - 8w_{30}\,w_{31} - 4w_{31}^{2} + w_{32} \\[2pt]
S_{19}={}& -19w_{8} + 8w_{8}^{2} + 8w_{8}\,w_{19} + 8w_{8}\,w_{23} - 46w_{9} - 24w_{9}^{2} + 24w_{9}\,w_{19} \\*
&\quad  + 24w_{9}\,w_{23} - 2w_{10} + 4w_{19} - 4w_{19}^{2} - 8w_{19}\,w_{23} + 14w_{23} - 4w_{23}^{2} - w_{25} \\*
&\quad  + w_{33} \\[2pt]
S_{20}={}& -3 - 6w_{8} - 4w_{10} - 2w_{25} - 4w_{33} - 4w_{33}^{2} + w_{34} \\[2pt]
S_{21}={}& 8 + 42w_{8} + 7w_{8}^{2} + 28w_{8}\,w_{20} + 48w_{9} - 32w_{9}^{2} - 32w_{9}\,w_{19} \\*
&\quad  + 64w_{9}\,w_{20} - 40w_{19}^{2} + 32w_{19}\,w_{20} + 14w_{20} - 4w_{20}^{2} + w_{35} + 2w_{38} \\[2pt]
S_{22}={}& 1 + 6w_{8} - 36w_{8}^{2} - 24w_{8}\,w_{35} + 2w_{35} - 4w_{35}^{2} + w_{36} \\[2pt]
S_{23}={}& -w_{9}^{2} - 4w_{9}\,w_{36} - 4w_{36}^{2} + w_{37} \\[2pt]
S_{24}={}& -6 - 2w_{8} - 4w_{9} - 14w_{19} - 16w_{19}^{2} + w_{38} \\[2pt]
S_{25}={}& 24 + 21w_{9} - 21w_{9}^{2} - 18w_{9}\,w_{10} - 24w_{9}\,w_{23} - 24w_{9}\,w_{38} + 3w_{10} \\*
&\quad  - 9w_{10}^{2} - 16w_{10}\,w_{38} + 3w_{23} - 12w_{23}^{2} - 8w_{23}\,w_{38} + 4w_{38} - 4w_{38}^{2} + w_{39} \\*
&\quad  + 2w_{46} \\[2pt]
S_{26}={}& 16 - 10w_{8} - 20w_{8}^{2} - 16w_{8}\,w_{10} - 24w_{8}\,w_{25} - 24w_{8}\,w_{39} - 12w_{10} \\*
&\quad  - 16w_{10}\,w_{25} - 16w_{10}\,w_{39} - 14w_{25} - 4w_{25}^{2} - 8w_{25}\,w_{39} + 4w_{39} - 4w_{39}^{2} \\*
&\quad  + w_{40} \\[2pt]
S_{27}={}& -w_{10}^{2} - 2w_{10}\,w_{11} - w_{11}^{2} + w_{41} \\[2pt]
S_{28}={}& 2 + w_{8}^{2} - 4w_{8}\,w_{41} + 2w_{10} - 2w_{10}^{2} - 8w_{10}\,w_{41} - 2w_{25} + 4w_{41} - 4w_{41}^{2} \\*
&\quad  + w_{42} \\[2pt]
S_{29}={}& -31 - 16w_{10} - 16w_{10}\,w_{23} + 16w_{10}\,w_{25} - 8w_{23} + 12w_{23}^{2} - 8w_{23}\,w_{25} + 24w_{25} \\*
&\quad  - 4w_{25}^{2} + w_{43} - 4w_{46} \\[2pt]
S_{30}={}& -4w_{10} - 16w_{25}^{2} - 32w_{25}\,w_{41} - 16w_{25}\,w_{42} - 16w_{25}\,w_{43} - 4w_{41} - 16w_{41}^{2} \\*
&\quad  - 16w_{41}\,w_{42} - 16w_{41}\,w_{43} - 2w_{42} - 4w_{42}^{2} - 8w_{42}\,w_{43} - 4w_{43}^{2} + w_{44} \\[2pt]
S_{31}={}& -16w_{8}^{2} + 8w_{8}\,w_{10} + 32w_{8}\,w_{25} + 16w_{8}\,w_{28} + w_{10} - w_{10}^{2} - 8w_{10}\,w_{25} \\*
&\quad  - 4w_{10}\,w_{28} - 16w_{25}^{2} - 16w_{25}\,w_{28} - 4w_{28}^{2} + w_{45} \\[2pt]
S_{32}={}& -1 - 8w_{10} - 4w_{10}^{2} + 8w_{10}\,w_{23} + 8w_{23} - 4w_{23}^{2} + w_{46} \\[2pt]
S_{33}={}& -1 - 9w_{10}^{2} + 24w_{10}\,w_{23} + 12w_{10}\,w_{46} - 16w_{23}^{2} - 16w_{23}\,w_{46} - 4w_{46}^{2} + w_{47} \\[2pt]
S_{34}={}& 1 + w_{9} - w_{9}^{2} - 4w_{9}\,w_{10} - 4w_{9}\,w_{47} - 4w_{10}^{2} - 8w_{10}\,w_{47} - 4w_{47}^{2} + w_{48} \\[2pt]
S_{35}={}& -1 - 4w_{48} - 4w_{48}^{2} + w_{49} \\[2pt]
S_{36}={}& -1 - w_{10}^{2} - 8w_{10}\,w_{48} - 4w_{10}\,w_{49} - 16w_{48}^{2} - 16w_{48}\,w_{49} - 4w_{49}^{2} + w_{50} \\[2pt]
S_{37}={}& 2 + w_{9} - w_{9}^{2} - 4w_{9}\,w_{10} - 4w_{9}\,w_{50} - 4w_{10}^{2} - 8w_{10}\,w_{50} - 4w_{50}^{2} + w_{51} \\[2pt]
S_{38}={}& -4w_{45}^{2} - 8w_{45}\,w_{51} - 4w_{51}^{2} + w_{52} \\[2pt]
S_{39}={}& -9 + 7w_{8} - w_{8}^{2} - 8w_{8}\,w_{45} - 4w_{8}\,w_{52} + 24w_{45} - 16w_{45}^{2} - 16w_{45}\,w_{52} \\*
&\quad  + 12w_{52} - 4w_{52}^{2} + w_{53} \\[2pt]
S_{40}={}& -1 - 4w_{53} - 4w_{53}^{2} + w_{54} \\[2pt]
S_{41}={}& -w_{10}^{2} - 8w_{10}\,w_{53} - 4w_{10}\,w_{54} - 16w_{53}^{2} - 16w_{53}\,w_{54} - 4w_{54}^{2} + w_{55} \\[2pt]
S_{42}={}& -w_{8} - w_{8}^{2} - 4w_{8}\,w_{10} - 4w_{8}\,w_{55} - 4w_{10} - 4w_{10}^{2} - 8w_{10}\,w_{55} - 4w_{55} \\*
&\quad  - 4w_{55}^{2} + w_{56} \\[2pt]
S_{43}={}& -10 - 12w_{56} - 4w_{56}^{2} + w_{57} \\[2pt]
S_{44}={}& 2 + 4w_{10} + 16w_{10}\,w_{56} + 8w_{10}\,w_{57} - 4w_{23} - 16w_{23}\,w_{56} - 8w_{23}\,w_{57} - w_{46} \\*
&\quad  + 8w_{56} - 16w_{56}^{2} - 16w_{56}\,w_{57} + 4w_{57} - 4w_{57}^{2} + w_{58} \\[2pt]
S_{45}={}& w_{9} - w_{9}^{2} + 8w_{9}\,w_{10} - 8w_{9}\,w_{23} - 4w_{9}\,w_{58} - 16w_{10}^{2} + 32w_{10}\,w_{23} \\*
&\quad  + 16w_{10}\,w_{58} - 16w_{23}^{2} - 16w_{23}\,w_{58} - 4w_{58}^{2} + w_{59} \\[2pt]
S_{46}={}& -10 - 12w_{51} - 4w_{51}^{2} + w_{60} \\[2pt]
S_{47}={}& 1 - 16w_{51}^{2} - 16w_{51}\,w_{59} - 16w_{51}\,w_{60} - 4w_{59}^{2} - 8w_{59}\,w_{60} - 4w_{60}^{2} + w_{61} \\[2pt]
S_{48}={}& w_{10} - w_{10}^{2} - 8w_{10}\,w_{59} - 4w_{10}\,w_{61} - 16w_{59}^{2} - 16w_{59}\,w_{61} - 4w_{61}^{2} + w_{62} \\[2pt]
S_{49}={}& -6w_{8} - 3w_{8}^{2} - 18w_{8}\,w_{9} - 16w_{9} - 27w_{9}^{2} + 4w_{9}\,w_{10} + 3w_{10} + 3w_{10}^{2} \\*
&\quad  - 4w_{10}\,w_{19} - 19w_{19} - 8w_{19}^{2} + w_{20} + w_{63} - 4w_{76} \\[2pt]
S_{50}={}& 1 - 8w_{8} - 4w_{8}\,w_{10} + 8w_{8}\,w_{63} - 16w_{9} - 8w_{9}\,w_{10} + 16w_{9}\,w_{63} + w_{10} \\*
&\quad  + w_{10}^{2} + 4w_{10}\,w_{19} + 8w_{19} - 8w_{19}\,w_{63} + 3w_{63} - 4w_{63}^{2} + w_{64} \\[2pt]
S_{51}={}& -2 - 12w_{8} - 16w_{8}\,w_{10} + 16w_{8}\,w_{62} + 16w_{8}\,w_{64} - 24w_{9} - 32w_{9}\,w_{10} \\*
&\quad  + 32w_{9}\,w_{62} + 32w_{9}\,w_{64} + 4w_{10} + 4w_{10}^{2} + 16w_{10}\,w_{19} + 12w_{19} - 16w_{19}\,w_{62} \\*
&\quad  - 16w_{19}\,w_{64} - 4w_{62}^{2} - 8w_{62}\,w_{64} - 4w_{63} - 2w_{64} - 4w_{64}^{2} + w_{65} \\[2pt]
S_{52}={}& -2 - 4w_{65} - 4w_{65}^{2} + w_{66} \\[2pt]
S_{53}={}& -w_{13} - 2w_{13}\,w_{227} + w_{67} \\[2pt]
S_{54}={}& w_{12} - w_{13} - 2w_{22}\,w_{67} - w_{68} \\[2pt]
S_{55}={}& w_{12} - w_{13} - 4w_{24}\,w_{67} - 2w_{27}\,w_{67} - w_{69} \\[2pt]
S_{56}={}& -2w_{29}\,w_{67} + 2w_{44}\,w_{67} - 2w_{67} - w_{70} + w_{75} \\[2pt]
S_{57}={}& w_{12} - w_{13} - 4w_{30}\,w_{67} - 2w_{32}\,w_{67} - w_{71} \\[2pt]
S_{58}={}& w_{12} - w_{13} - 2w_{34}\,w_{67} + 6w_{67} - w_{72} \\[2pt]
S_{59}={}& w_{12} - w_{13} - 4w_{36}\,w_{67} - 2w_{37}\,w_{67} - w_{73} \\[2pt]
S_{60}={}& w_{12} - w_{13} - 2w_{40}\,w_{67} - w_{74} \\[2pt]
S_{61}={}& w_{12} - w_{13} - 2w_{44}\,w_{67} - 4w_{67} - w_{75} \\[2pt]
S_{62}={}& -w_{8}\,w_{9} + w_{8}\,w_{10} + w_{9} + w_{9}\,w_{10} - w_{10} - w_{10}^{2} + w_{76} \\[2pt]
S_{63}={}& 1 + 2w_{11} + w_{11}^{2} - w_{76}^{2} + w_{77} \\[2pt]
S_{64}={}& -w_{77}^{2} + w_{78} \\[2pt]
S_{65}={}& -4w_{65}\,w_{67} - 2w_{66}\,w_{67} + 2w_{67} + w_{79} \\[2pt]
S_{66}={}& -w_{12}^{2} + 2w_{12}\,w_{13} + 2w_{12}\,w_{79} - w_{13}^{2} - 2w_{13}\,w_{79} - w_{79}^{2} + w_{80} \\[2pt]
S_{67}={}& -1 - w_{9} - 2w_{9}\,w_{227} + w_{81} \\[2pt]
S_{68}={}& -w_{13}\,w_{81} + w_{83}\,w_{84} + w_{85} \\[2pt]
S_{69}={}& -w_{8} - 2w_{8}\,w_{227} + w_{82} + 2w_{227} \\[2pt]
S_{70}={}& 1 + 2w_{8} + 4w_{8}\,w_{227} + 4w_{9} + 8w_{9}\,w_{227} - 2w_{19} - 4w_{19}\,w_{227} + w_{83} + 4w_{227} \\[2pt]
S_{71}={}& -w_{12}\,w_{82} + w_{84} \\[2pt]
S_{72}={}& w_{12}\,w_{81} + 6w_{13} - 8w_{13}\,w_{81} - 4w_{13}\,w_{82} - w_{13}\,w_{83} - 12w_{13}\,w_{227} + w_{86} \\[2pt]
S_{73}={}& -w_{12}\,w_{68} + w_{13}\,w_{68} + w_{87} \\[2pt]
S_{74}={}& -w_{69}\,w_{87} + w_{88} \\[2pt]
S_{75}={}& -w_{70}\,w_{88} + w_{89} \\[2pt]
S_{76}={}& -w_{71}\,w_{89} + w_{90} \\[2pt]
S_{77}={}& -w_{72}\,w_{90} + w_{91} \\[2pt]
S_{78}={}& -w_{73}\,w_{91} + w_{92} \\[2pt]
S_{79}={}& -w_{74}\,w_{92} + w_{93} \\[2pt]
S_{80}={}& -w_{75}\,w_{93} + w_{94} \\[2pt]
S_{81}={}& -w_{85}\,w_{94} + w_{95} \\[2pt]
S_{82}={}& -w_{86}\,w_{95} + w_{96} \\[2pt]
S_{83}={}& -w_{96}^{2} + w_{97} \\[2pt]
S_{84}={}& -w_{2} - w_{2}\,w_{16} + w_{16} + w_{98} \\[2pt]
S_{85}={}& -w_{14}\,w_{98} + w_{99} \\[2pt]
S_{86}={}& -2w_{1}^{2} + 2w_{1}\,w_{12} + 2w_{1}\,w_{13} + 2w_{1}\,w_{99} + 2w_{1}\,w_{100} - w_{12}^{2} - 2w_{12}\,w_{99} \\*
&\quad  - w_{13}^{2} - 2w_{13}\,w_{100} + w_{17}\,w_{101} - w_{78}\,w_{97} - w_{80}\,w_{97} - w_{99}^{2} - w_{100}^{2} \\[2pt]
S_{87}={}& -w_{2}\,w_{15} + w_{15} - w_{15}\,w_{98} + w_{100} \\[2pt]
S_{88}={}& -3w_{104} - 6w_{104}\,w_{220} + w_{106} \\[2pt]
S_{89}={}& w_{102}\,w_{211} - w_{105} - w_{106}\,w_{127} - w_{127} - 2w_{127}\,w_{227} \\[2pt]
S_{90}={}& w_{107}\,w_{212} - w_{109}\,w_{130} - w_{110} \\[2pt]
S_{91}={}& w_{111}\,w_{212} - w_{112}\,w_{130} - w_{113} \\[2pt]
S_{92}={}& w_{114}\,w_{212} - w_{115}\,w_{130} - w_{116} \\[2pt]
S_{93}={}& w_{117}\,w_{212} - w_{118}\,w_{130} - w_{119} \\[2pt]
S_{94}={}& w_{120}\,w_{213} - w_{122}\,w_{130} - w_{123} \\[2pt]
S_{95}={}& w_{124}\,w_{213} - w_{125}\,w_{130} - w_{126} \\[2pt]
S_{96}={}& -w_{17}\,w_{127} + w_{130} \\[2pt]
S_{97}={}& -w_{130}\,w_{134} + w_{133}\,w_{211} - w_{135} \\[2pt]
S_{98}={}& -w_{130}\,w_{137} + w_{136}\,w_{212} - w_{138} \\[2pt]
S_{99}={}& w_{139} - 2w_{139}\,w_{141} + w_{142} \\[2pt]
S_{100}={}& -w_{139}\,w_{142} - w_{142} + w_{143} \\[2pt]
S_{101}={}& -1 + w_{142}^{2} - w_{143}^{2} + w_{223}^{2} \\[2pt]
S_{102}={}& -w_{140}\,w_{145} - 7w_{145} + w_{146} \\[2pt]
S_{103}={}& 2w_{144} + w_{144}^{2} + w_{145}^{2} - w_{146}^{2} \\[2pt]
S_{104}={}& 1 - w_{4} - 32w_{140}\,w_{147} + w_{144} + 16w_{145} - w_{146} + 33w_{147} \\[2pt]
S_{105}={}& -12 - 2w_{140} - 4w_{140}\,w_{227} + w_{148} - 24w_{227} + 4w_{227}^{2} \\[2pt]
S_{106}={}& w_{17} + w_{148}\,w_{149} - w_{182} - w_{183} - 2w_{183}\,w_{227} + w_{184} \\[2pt]
S_{107}={}& -1 + 2w_{140}\,w_{195} - w_{150} + 14w_{195} - w_{195}^{2} \\[2pt]
S_{108}={}& 3 + w_{150}\,w_{151} - w_{182} - w_{183}\,w_{195} + w_{184} + 6w_{220} \\[2pt]
S_{109}={}& -w_{140}\,w_{153} - 7w_{153} + w_{154} \\[2pt]
S_{110}={}& 2w_{152} + w_{152}^{2} + w_{153}^{2} - w_{154}^{2} \\[2pt]
S_{111}={}& -39 + w_{130}\,w_{228} - 8w_{140} - 12w_{140}\,w_{220} - 36w_{220} + 36w_{220}^{2} \\[2pt]
S_{112}={}& 1 - w_{102} - w_{140}\,w_{153} + w_{152} - 3w_{153} + 6w_{153}\,w_{220} \\[2pt]
S_{113}={}& 7 - w_{130}\,w_{228} + 2w_{140} + w_{155} - 12w_{220} \\[2pt]
S_{114}={}& -1 + w_{103} - w_{144} - 3w_{145} - 6w_{145}\,w_{220} + w_{146} + w_{155}\,w_{156} \\[2pt]
S_{115}={}& 7 + w_{140} - w_{140}\,w_{158} - 7w_{158} + w_{159} \\[2pt]
S_{116}={}& 1 + 2w_{157} + w_{157}^{2} - 2w_{158} + w_{158}^{2} - w_{159}^{2} \\[2pt]
S_{117}={}& 4 - w_{107} + w_{140} - w_{140}\,w_{158} + w_{157} - 3w_{158} + 6w_{158}\,w_{220} - 6w_{220} \\[2pt]
S_{118}={}& -3w_{103} - 6w_{103}\,w_{220} + w_{108} \\[2pt]
S_{119}={}& -w_{140}\,w_{161} - 7w_{161} + w_{162} \\[2pt]
S_{120}={}& 2w_{160} + w_{160}^{2} + w_{161}^{2} - w_{162}^{2} \\[2pt]
S_{121}={}& 1 - w_{111} - w_{140}\,w_{161} + w_{160} - 3w_{161} + 6w_{161}\,w_{220} \\[2pt]
S_{122}={}& -w_{140}\,w_{164} - 7w_{164} + w_{165} \\[2pt]
S_{123}={}& 2w_{163} + w_{163}^{2} + w_{164}^{2} - w_{165}^{2} \\[2pt]
S_{124}={}& 1 - w_{114} - w_{140}\,w_{164} + w_{163} - 3w_{164} + 6w_{164}\,w_{220} \\[2pt]
S_{125}={}& -w_{140}\,w_{167} - 7w_{167} + w_{168} \\[2pt]
S_{126}={}& 2w_{166} + w_{166}^{2} + w_{167}^{2} - w_{168}^{2} \\[2pt]
S_{127}={}& 1 - w_{117} - w_{140}\,w_{167} + w_{166} - 3w_{167} + 6w_{167}\,w_{220} \\[2pt]
S_{128}={}& -w_{140}\,w_{170} - 7w_{170} + w_{171} \\[2pt]
S_{129}={}& 2w_{169} + w_{169}^{2} + w_{170}^{2} - w_{171}^{2} \\[2pt]
S_{130}={}& 1 - w_{120} - w_{140}\,w_{170} + w_{169} - 3w_{170} + 6w_{170}\,w_{220} \\[2pt]
S_{131}={}& -w_{140}\,w_{173} - 7w_{173} + w_{174} \\[2pt]
S_{132}={}& 2w_{172} + w_{172}^{2} + w_{173}^{2} - w_{174}^{2} \\[2pt]
S_{133}={}& -1 + w_{121} + w_{155}\,w_{175} - w_{172} - 3w_{173} - 6w_{173}\,w_{220} + w_{174} \\[2pt]
S_{134}={}& -w_{140}\,w_{177} - 7w_{177} + w_{178} \\[2pt]
S_{135}={}& 2w_{176} + w_{176}^{2} + w_{177}^{2} - w_{178}^{2} \\[2pt]
S_{136}={}& 1 - w_{124} - w_{140}\,w_{177} + w_{176} - 3w_{177} + 6w_{177}\,w_{220} \\[2pt]
S_{137}={}& -w_{140}\,w_{180} - 7w_{180} + w_{181} \\[2pt]
S_{138}={}& 2w_{179} + w_{179}^{2} + w_{180}^{2} - w_{181}^{2} \\[2pt]
S_{139}={}& 1 - w_{128}\,w_{130} - w_{129} - w_{140}\,w_{180} + w_{179} - 3w_{180} + 6w_{180}\,w_{220} \\[2pt]
S_{140}={}& 1 - w_{140}\,w_{183} - 7w_{183} + w_{184} \\[2pt]
S_{141}={}& -1 + 2w_{182} + w_{182}^{2} + w_{183}^{2} - 2w_{184} - w_{184}^{2} \\[2pt]
S_{142}={}& w_{127} + w_{155}\,w_{185} - w_{182} - 3w_{183} - 6w_{183}\,w_{220} + w_{184} \\[2pt]
S_{143}={}& -w_{140}\,w_{187} - 7w_{187} + w_{188} \\[2pt]
S_{144}={}& 2w_{186} + w_{186}^{2} + w_{187}^{2} - w_{188}^{2} \\[2pt]
S_{145}={}& 1 - w_{130}\,w_{131} - w_{132} - w_{140}\,w_{187} + w_{186} - 3w_{187} + 6w_{187}\,w_{220} \\[2pt]
S_{146}={}& -w_{140}\,w_{190} - 7w_{190} + w_{191} \\[2pt]
S_{147}={}& 2w_{189} + w_{189}^{2} + w_{190}^{2} - w_{191}^{2} \\[2pt]
S_{148}={}& 1 - w_{133} - w_{140}\,w_{190} + w_{189} - 3w_{190} + 6w_{190}\,w_{220} \\[2pt]
S_{149}={}& -w_{140}\,w_{193} - 7w_{193} + w_{194} \\[2pt]
S_{150}={}& 2w_{192} + w_{192}^{2} + w_{193}^{2} - w_{194}^{2} \\[2pt]
S_{151}={}& 1 - w_{136} - w_{140}\,w_{193} + w_{192} - 3w_{193} + 6w_{193}\,w_{220} \\[2pt]
S_{152}={}& -10 - 2w_{0}^{2} - 2w_{0}\,w_{4} - w_{1}^{2} - w_{2}^{2} - w_{4}^{2} + 2w_{8} + 2w_{8}\,w_{9} - w_{10}^{2} \\*
&\quad  - w_{11}^{2} - w_{12}^{2} - w_{13}^{2} - w_{14}^{2} - w_{15}^{2} - w_{16}^{2} - w_{17}^{2} - w_{19} - w_{103}^{2} - w_{121}^{2} \\*
&\quad  - w_{127}^{2} + w_{139} + w_{145}^{2} + w_{153}^{2} - 2w_{158} + w_{158}^{2} + w_{161}^{2} + w_{164}^{2} + w_{167}^{2} + w_{170}^{2} \\*
&\quad  + w_{173}^{2} + w_{177}^{2} + w_{180}^{2} + w_{183}^{2} + w_{187}^{2} + w_{190}^{2} - w_{193}^{2} - w_{195}^{2} + w_{214}^{2} + w_{215}^{2} \\*
&\quad  + w_{216} + w_{216}^{2} - 36w_{220} - 36w_{220}^{2} - 4w_{227} - 4w_{227}^{2} \\[2pt]
S_{153}={}& w_{195} - w_{226} - 2w_{226}\,w_{227} \\[2pt]
S_{154}={}& -16 + w_{0} - w_{196}^{2} \\[2pt]
S_{155}={}& -w_{140}\,w_{198} - 7w_{198} + w_{207} \\[2pt]
S_{156}={}& -1 + w_{197}^{2} + w_{198}^{2} - w_{207}^{2} \\[2pt]
S_{157}={}& -w_{200} - w_{200}\,w_{201} + w_{208} \\[2pt]
S_{158}={}& -1 + w_{199}^{2} + w_{200}^{2} - w_{208}^{2} \\[2pt]
S_{159}={}& -w_{193}^{2} + w_{209} \\[2pt]
S_{160}={}& w_{198} - 2w_{203}\,w_{209} - w_{209} \\[2pt]
S_{161}={}& -4w_{198} - 8w_{198}\,w_{202} + w_{201} \\[2pt]
S_{162}={}& -6 - w_{140} - w_{197}\,w_{204} + w_{201} \\[2pt]
S_{163}={}& -1 - w_{192} - w_{197}\,w_{205} + w_{199} \\[2pt]
S_{164}={}& -w_{8} - w_{9} - 4w_{19} - 2w_{22} - 4w_{24} - 2w_{27} - 2w_{29} - 4w_{30} - 2w_{32} - 2w_{34} \\*
&\quad  - 4w_{36} - 2w_{37} - 2w_{40} - 2w_{44} - 4w_{65} - 2w_{66} - 4w_{193}\,w_{206} + w_{200} \\[2pt]
S_{165}={}& -2 + 2w_{0}^{2} + 2w_{0}\,w_{4} - 2w_{0}\,w_{145} - 2w_{0}\,w_{153} + w_{1}^{2} - 2w_{1}\,w_{183} + w_{2}^{2} \\*
&\quad  - 2w_{2}\,w_{173} + w_{4}^{2} - 2w_{4}\,w_{153} + w_{8}^{2} - 2w_{8}\,w_{158} - 2w_{8}\,w_{193} + w_{9}^{2} \\*
&\quad  - 2w_{9}\,w_{161} - 2w_{9}\,w_{193} + w_{10}^{2} - 2w_{10}\,w_{164} + w_{11}^{2} - 2w_{11}\,w_{167} + w_{12}^{2} \\*
&\quad  - 2w_{12}\,w_{170} + w_{13}^{2} - 2w_{13}\,w_{177} + w_{14}^{2} - 2w_{14}\,w_{180} + w_{15}^{2} - 2w_{15}\,w_{187} + w_{16}^{2} \\*
&\quad  - 2w_{16}\,w_{190} - 8w_{19}\,w_{193} - 4w_{22}\,w_{193} - 8w_{24}\,w_{193} - 4w_{27}\,w_{193} - 4w_{29}\,w_{193} \\*
&\quad  - 8w_{30}\,w_{193} - 4w_{32}\,w_{193} - 4w_{34}\,w_{193} - 8w_{36}\,w_{193} - 4w_{37}\,w_{193} - 4w_{40}\,w_{193} \\*
&\quad  - 4w_{44}\,w_{193} - 8w_{65}\,w_{193} - 4w_{66}\,w_{193} - w_{140}\,w_{210} + w_{144}^{2} + w_{152}^{2} + w_{157}^{2} + 2w_{158} \\*
&\quad  + w_{160}^{2} + w_{163}^{2} + w_{166}^{2} + w_{169}^{2} + w_{172}^{2} + w_{176}^{2} + w_{179}^{2} + w_{182}^{2} + w_{186}^{2} + w_{189}^{2} \\*
&\quad  + w_{192}^{2} + 2w_{193}^{2} - 6w_{210} - w_{214}^{2} - w_{215}^{2} - w_{216} - w_{216}^{2} \\[2pt]
S_{166}={}& -w_{103}\,w_{211} + w_{127} \\[2pt]
S_{167}={}& -w_{108}\,w_{212} + w_{127} \\[2pt]
S_{168}={}& -w_{121}\,w_{213} + w_{127} \\[2pt]
S_{169}={}& -1 + w_{2} - w_{3}\,w_{221} \\[2pt]
S_{170}={}& -1 - 8w_{8}\,w_{19} - 4w_{8}\,w_{22} - 8w_{8}\,w_{24} - 4w_{8}\,w_{27} - 4w_{8}\,w_{29} - 8w_{8}\,w_{30} \\*
&\quad  - 4w_{8}\,w_{32} - 4w_{8}\,w_{34} - 8w_{8}\,w_{36} - 4w_{8}\,w_{37} - 4w_{8}\,w_{40} - 4w_{8}\,w_{44} \\*
&\quad  - 8w_{8}\,w_{65} - 4w_{8}\,w_{66} - 8w_{9}\,w_{19} - 4w_{9}\,w_{22} - 8w_{9}\,w_{24} - 4w_{9}\,w_{27} \\*
&\quad  - 4w_{9}\,w_{29} - 8w_{9}\,w_{30} - 4w_{9}\,w_{32} - 4w_{9}\,w_{34} - 8w_{9}\,w_{36} - 4w_{9}\,w_{37} \\*
&\quad  - 4w_{9}\,w_{40} - 4w_{9}\,w_{44} - 8w_{9}\,w_{65} - 4w_{9}\,w_{66} - w_{19} - 16w_{19}^{2} - 16w_{19}\,w_{22} \\*
&\quad  - 32w_{19}\,w_{24} - 16w_{19}\,w_{27} - 16w_{19}\,w_{29} - 32w_{19}\,w_{30} - 16w_{19}\,w_{32} - 16w_{19}\,w_{34} \\*
&\quad  - 32w_{19}\,w_{36} - 16w_{19}\,w_{37} - 16w_{19}\,w_{40} - 16w_{19}\,w_{44} - 32w_{19}\,w_{65} - 16w_{19}\,w_{66} \\*
&\quad  - 4w_{22}^{2} - 16w_{22}\,w_{24} - 8w_{22}\,w_{27} - 8w_{22}\,w_{29} - 16w_{22}\,w_{30} - 8w_{22}\,w_{32} \\*
&\quad  - 8w_{22}\,w_{34} - 16w_{22}\,w_{36} - 8w_{22}\,w_{37} - 8w_{22}\,w_{40} - 8w_{22}\,w_{44} - 16w_{22}\,w_{65} \\*
&\quad  - 8w_{22}\,w_{66} - 16w_{24}^{2} - 16w_{24}\,w_{27} - 16w_{24}\,w_{29} - 32w_{24}\,w_{30} - 16w_{24}\,w_{32} \\*
&\quad  - 16w_{24}\,w_{34} - 32w_{24}\,w_{36} - 16w_{24}\,w_{37} - 16w_{24}\,w_{40} - 16w_{24}\,w_{44} - 32w_{24}\,w_{65} \\*
&\quad  - 16w_{24}\,w_{66} - 4w_{27}^{2} - 8w_{27}\,w_{29} - 16w_{27}\,w_{30} - 8w_{27}\,w_{32} - 8w_{27}\,w_{34} \\*
&\quad  - 16w_{27}\,w_{36} - 8w_{27}\,w_{37} - 8w_{27}\,w_{40} - 8w_{27}\,w_{44} - 16w_{27}\,w_{65} - 8w_{27}\,w_{66} \\*
&\quad  - 4w_{29}^{2} - 16w_{29}\,w_{30} - 8w_{29}\,w_{32} - 8w_{29}\,w_{34} - 16w_{29}\,w_{36} - 8w_{29}\,w_{37} \\*
&\quad  - 8w_{29}\,w_{40} - 8w_{29}\,w_{44} - 16w_{29}\,w_{65} - 8w_{29}\,w_{66} - 16w_{30}^{2} - 16w_{30}\,w_{32} \\*
&\quad  - 16w_{30}\,w_{34} - 32w_{30}\,w_{36} - 16w_{30}\,w_{37} - 16w_{30}\,w_{40} - 16w_{30}\,w_{44} - 32w_{30}\,w_{65} \\*
&\quad  - 16w_{30}\,w_{66} - 4w_{32}^{2} - 8w_{32}\,w_{34} - 16w_{32}\,w_{36} - 8w_{32}\,w_{37} - 8w_{32}\,w_{40} \\*
&\quad  - 8w_{32}\,w_{44} - 16w_{32}\,w_{65} - 8w_{32}\,w_{66} - 4w_{34}^{2} - 16w_{34}\,w_{36} - 8w_{34}\,w_{37} \\*
&\quad  - 8w_{34}\,w_{40} - 8w_{34}\,w_{44} - 16w_{34}\,w_{65} - 8w_{34}\,w_{66} - 16w_{36}^{2} - 16w_{36}\,w_{37} \\*
&\quad  - 16w_{36}\,w_{40} - 16w_{36}\,w_{44} - 32w_{36}\,w_{65} - 16w_{36}\,w_{66} - 4w_{37}^{2} - 8w_{37}\,w_{40} \\*
&\quad  - 8w_{37}\,w_{44} - 16w_{37}\,w_{65} - 8w_{37}\,w_{66} - 4w_{40}^{2} - 8w_{40}\,w_{44} - 16w_{40}\,w_{65} \\*
&\quad  - 8w_{40}\,w_{66} - 4w_{44}^{2} - 16w_{44}\,w_{65} - 8w_{44}\,w_{66} - 16w_{65}^{2} - 16w_{65}\,w_{66} - 4w_{66}^{2} \\*
&\quad  - w_{145}^{2} - w_{153}^{2} + 2w_{158} - w_{158}^{2} - w_{161}^{2} - w_{164}^{2} - w_{167}^{2} - w_{170}^{2} - w_{173}^{2} - w_{177}^{2} \\*
&\quad  - w_{180}^{2} - w_{183}^{2} - w_{187}^{2} - w_{190}^{2} + w_{193}^{2} - w_{214}^{2} - w_{215}^{2} - w_{216} - w_{216}^{2} \\[2pt]
S_{171}={}& -1 - 4w_{105}^{2} - 4w_{110}^{2} - 4w_{113}^{2} - 4w_{116}^{2} - 4w_{119}^{2} - 4w_{123}^{2} - 4w_{126}^{2} + w_{127}^{2} \\*
&\quad  - 4w_{129}^{2} - 4w_{132}^{2} - 4w_{135}^{2} - 4w_{138}^{2} - w_{217}^{2} - w_{218}^{2} - w_{219} - w_{219}^{2} \\[2pt]
S_{172}={}& -1 + w_{221} - 4w_{222} - 4w_{222}^{2} \\[2pt]
S_{173}={}& 7 + w_{140} - w_{223}^{2} - w_{224}^{2} - w_{225}^{2} \\[2pt]
S_{174}={}& -3 - 2w_{0} - 2w_{4} + 2w_{8} - w_{8}^{2} - w_{9}^{2} - w_{10}^{2} - w_{11}^{2} - w_{12}^{2} - w_{13}^{2} \\*
&\quad  - w_{14}^{2} - w_{15}^{2} - w_{16}^{2} + w_{226} \\
\end{align*}
\endgroup
% END report/compact_residuals.tex
\end{document}
